\documentclass[11pt]{amsart}
\usepackage{amsmath,amssymb,amsthm}
\usepackage[margin=1.1in]{geometry}
\usepackage{hyperref}

\theoremstyle{plain}
\newtheorem{theorem}{Theorem}[section]
\newtheorem{lemma}[theorem]{Lemma}
\newtheorem{proposition}[theorem]{Proposition}
\newtheorem{corollary}[theorem]{Corollary}
\theoremstyle{definition}
\newtheorem{definition}[theorem]{Definition}
\newtheorem{remark}[theorem]{Remark}
\newtheorem{conjecture}[theorem]{Conjecture}

\newcommand{\LR}{\operatorname{LR}}
\newcommand{\nrm}[1]{\left\lVert #1 \right\rVert}
\newcommand{\Z}{\mathbb{Z}}
\newcommand{\R}{\mathbb{R}}

\begin{document}

\title{Two- and three-speed modifications of the tight Lonely Runner instance:\\
exhaustive verification, and a Lean~4 formalization of the covering framework}
\author{Vico Bonfioli}
\email{vicobonfioli@gmail.com}

\begin{abstract}
For a set $V$ of $n-1$ distinct positive integers write $\LR(V)=\max_t \min_{v\in V}\nrm{vt}$, where
$\nrm{x}$ is the distance from $x$ to the nearest integer; $V$ is \emph{tight} if $\LR(V)=1/n$, the
value predicted by the Lonely Runner Conjecture. Every such $V$ is a \emph{$k$-speed modification}
of the baseline $[n-1]=\{1,\dots,n-1\}$: writing $M=[n-1]\setminus V$ and $R=V\setminus[n-1]$ one
has $|M|=|R|=k$, where $k$ counts the speeds of $V$ that are at least $n$.

The case $k=1$ is settled. Goddyn and Wong \cite{gw} classified the tight sets obtained by replacing
one speed with a multiple, and Zhang \cite{zhang} (version~2, 31 August 2026) completed the
classification of one-entry modifications (Thm.~1.2), bounded the largest speed by $O(n\log n)$
(Thm.~1.3), and computed exactly the components of the region $U(n,r)$ left uncovered when a single
speed $r$ is deleted (Thm.~1.4). For $k=2$, Zhang \cite[Thms.~7.5, 7.10]{zhang} proves that no tight
two-speed modification exists for removed sets $\{2,r'\}$ with $n\ge19$ and $\{q,r'\}$ with
$n\ge\max(40,10q)$, with no bound on the inserted speeds, and his census \cite[Thm.~7.13]{zhang}
covers $8\le n\le140$ with inserted speeds at most $8n$. Version~1 of that paper had verified
$6\le n\le34$ under the hypothesis that the inserted speeds are at most $4n$.

This paper contributes three things. First, a covering bound which is uniform over the inserted
speeds: if $\LR(V)\le1/n$ then $\sum_{v\in R}1/v\ \ge\ \ell(n,M)(n-2k)/2$, where $\ell(n,M)$ is the
length of the longest component of the uncovered region, with \emph{no} upper bound on any inserted
speed assumed. For $k=2$ Zhang's non-existence theorems \cite[Thms.~7.5, 7.10]{zhang} are stronger than the bound
here on the removed sets they cover; the measure argument here is uniform in $k$ and in the inserted
speeds, at the cost of a constant. Second, and as a
consequence, an exhaustive verification with the cap on the inserted speeds removed: the conjecture
holds for every $2$-speed modification with $5\le n\le100$ and every $3$-speed modification with
$7\le n\le40$. The enumeration is finite for two rigorous reasons and no branch was truncated.
Third, a formalization in Lean~4 with Mathlib of the covering framework and of the description of
the uncovered region, comprising $70$ theorems and $5$ definitions with no \texttt{sorry}. We are not aware of any
previous formalization of the material treated here; for the conjecture itself, the Acorn library
contains proofs of statements with the constant slack by one, discussed in the companion paper. The statement itself, and Tao's
bound, appear in Lean in the \texttt{formal-conjectures} repository with \texttt{sorry} bodies and
no proof content. For the wider context of the problem we refer to the survey \cite{ps}. Not every statement below is
formalized, and the Formalization section is the authoritative account of which are and which are
not; downstream of the strip lemma the formalized constants carry $2+4/n$ in place of the sharp $2$.

Much of Sections~\ref{sec:uncond}, \ref{sec:region}, \ref{sec:localize} and~\ref{sec:divis} is
expository; Section~\ref{sec:comp} and the formalization are not. The description of the uncovered region
by localization at fractions whose denominators were deleted is Zhang's \cite[\S2, proof of Lemma~2.4]{zhang}, used
constructively by Bohman and Peng \cite[\S3]{bp}; the measure of the uncovered region as an object of
study is Cusick's \cite{cusick74}, with measure and component-count bookkeeping already in
Cusick and Pomerance \cite{cp}; the gap-and-block framework is Bohman, Holzman and Kleitman
\cite{bhk}; the divisibility obstruction is \cite[Lemma~2.2]{gw}, \cite[Lemmas~2.3, 6.1]{zhang} and
\cite[Lemma~4]{rosenfeld}; and the factor-of-two barrier for measure arguments is \cite[Prop.~1.1]{tao}.
We include these because the formalization and the search consume them, and we claim no priority for
any of them.
\end{abstract}

Statements are tagged with the status of their proof: \emph{Verified} means formalized in Lean~4
against Mathlib with no \texttt{sorry} and no axioms beyond \texttt{propext},
\texttt{Classical.choice} and \texttt{Quot.sound}; \emph{Computed} means established by a finite
computation, with the generating script cited. An untagged theorem, lemma, proposition or corollary
carries a complete proof in the text and no formalization. Conjectures are labelled as such and are
proved nowhere here. Definitions carry no status.


\date{18 October 2026. Preprint, not peer reviewed.}
\maketitle

\section{Introduction}

Let $\nrm{x}$ be the distance from $x\in\R$ to the nearest integer. For a set $V$ of $n-1$ distinct
positive integers put $\LR(V)=\max_{t\in\R}\min_{v\in V}\nrm{vt}$. The Lonely Runner Conjecture
asserts $\LR(V)\ge 1/n$. The baseline $[n-1]=\{1,\dots,n-1\}$ satisfies $\LR([n-1])=1/n$, attained at
$t=1/n$, so it is \emph{tight}.

Goddyn and Wong \cite{gw} classified the tight sets obtained from $[n-1]$ by replacing one speed $r$
with a multiple, and proved that for fixed $r$ only finitely many non-multiple replacements can be
tight. Zhang \cite{zhang} made that finiteness effective. Writing $U(n,r)$ for the region left uncovered when $r$ is deleted (Definition~\ref{def:U} below,
with $M=\{r\}$), he computes the length of every connected component of $U(n,r)$ exactly, in
both regimes $2r>n-1$ and $2r\le n-1$, and deduces an explicit upper bound on the inserted speed, hence $n\le 6r$ when $2r\le n-1$,
together with the complete classification for $r=2$ and $r=3$.

Throughout, $P(b)$ denotes the smallest prime not dividing $b$.

This note treats modifications of arbitrary size.

\begin{definition}\label{def:mod}
Let $n\ge 3$ and $k\ge1$. A \emph{$k$-speed modification} of the baseline is
\[
  V\;=\;\bigl([n-1]\setminus M\bigr)\cup R,\qquad M\subseteq[n-1],\ |M|=|R|=k,\ \ v\ge n \text{ for all } v\in R .
\]
\end{definition}

\begin{remark}\label{rem:general}
This is no restriction on the speed set. Let $V$ be any set of $n-1$ distinct positive integers, and
put $M=[n-1]\setminus V$, $R=V\setminus[n-1]$, $k=\#\{v\in V:v\ge n\}$. Then $|V\cap[n-1]|=n-1-k$, so
$|M|=(n-1)-(n-1-k)=k=|R|$, and $V$ is the $k$-speed modification determined by $M$ and $R$. Every
instance with $1\le k\le n-2$ arises exactly once in this way, with $k$ recording how
many speeds exceed the baseline. The hypothesis $2k<n$ belongs to Theorem~\ref{thm:cover}, not to the
setting: Theorem~\ref{thm:localize} holds for every $k\le n-2$, and Theorem~\ref{thm:exact} for
every $k\le n-2$ as well.
\end{remark}

Then $|V|=n-1$, so the predicted floor for $V$ is again $1/n$. The inserted speeds are subject to no
upper bound, and that is the difficulty.

\begin{definition}\label{def:U}
For $M\subseteq[n-1]$ and $v\ge1$ put
\[
  U(n,M)=\bigl\{t\in[0,1):\nrm{jt}\ge 1/n \ \text{ for all } j\in[n-1]\setminus M\bigr\},
  \qquad
  B_v=\bigl\{t\in[0,1):\nrm{vt}<1/n\bigr\}.
\]
Both depend on $n$; the dependence of $B_v$ on $n$ is suppressed. Since $U(n,M)$ is a finite union
of closed intervals, a longest component exists; write $\ell(n,M)$ for its length, which is $0$
exactly when $U(n,M)$ is finite.
\end{definition}

$U(n,M)$ is closed and each $B_v$ is open of Lebesgue measure exactly $2/n$, independently of $v$.
Directly from the definitions, $\LR(V)\ge 1/n$ if and only if
$U(n,M)\not\subseteq\bigcup_{v\in R}B_v$.

\subsection*{What is new here, and what is not}

Arguments that show a family of bad sets fails to cover the uncovered region of a sub-family are
long established. Bohman, Holzman and Kleitman \cite{bhk} work with exactly this dichotomy, calling
a maximal uncovered interval for a sub-collection $S$ an $S$-gap and a maximal covered interval an
$S$-block; their Lemma~4 bounds the length of a block of the \emph{two} fastest speeds and their
Corollary~5 concludes non-covering from a long enough gap. So the passage from one covering speed to
two is not new, and neither is the idea of measuring the uncovered region, which goes back to
Cusick \cite{cusick74} and, with component counting alongside the measure, to Cusick and Pomerance
\cite{cp}.

What is new is the form of the bound. For a single inserted speed the natural step is a
connectivity argument: a connected subset of $\{t:\nrm{wt}\le 1/n\}$ lies inside a single one of its
constituent intervals, so its length is at most $2/(wn)$. That uses in an essential way that there
is only \emph{one} inserted speed. The argument of \cite{bhk} handles two, but through explicit
block-length estimates tied to the five- and six-runner case analysis. Section~7 of \cite{zhang}
(version~2) handles two inserted speeds by a counting bound on the part of a bad set inside a
component (Lemma~7.3) and proves the non-existence statements of Theorems~7.5 and~7.10 with no bound
on the inserted speeds; on the removed sets they cover, for $k=2$, these are stronger than
Theorem~\ref{thm:cover}, which only bounds $\sum_{v\in R}1/v$. Theorem~\ref{thm:cover} replaces
these arguments by a measure estimate which holds for every $k$ simultaneously, is uniform over the
inserted speeds, and assumes no upper bound on any of them. That uniformity is what makes the
enumeration of Section~\ref{sec:comp} finite for $k=3$ without a cap on $\max R$; for $k=2$ the
cap that version~1 of \cite{zhang} needed is removed there by other means.

The price is a constant. Theorem~\ref{thm:cover} carries a factor $2$ that the connectivity argument
does not, so at $k=1$ it is weaker than \cite[Thm.~1.4]{zhang}, which computes $\ell(n,\{b\})$
exactly. Sections~\ref{sec:uncond} to~\ref{sec:region} accordingly record the $k=1$ theory in the
form the later sections consume, with attribution; the reader wanting the sharp $k=1$ statements
should take them from \cite{zhang}.

\begin{theorem}[Proved; Covering bound]\label{thm:cover}
Let $V$ be a $k$-speed modification with $2k<n$ and $\LR(V)\le 1/n$. Then
\[
  \sum_{v\in R}\frac1v\;\ge\;\frac{\ell(n,M)\,(n-2k)}{2},
  \qquad\text{and consequently}\qquad
  \min R\;\le\;\frac{2k}{\ell(n,M)\,(n-2k)}
\]
whenever $\ell(n,M)>0$. No upper bound on any element of $R$ is assumed.
\end{theorem}

A large inserted speed is close to equidistributed and contributes about $2\ell/n$ to a component of
length $\ell$; $k$ of them contribute about $2k\ell/n<\ell$ and cannot finish the covering. The
deficit must come from the boundary term $2/(vn)$, which is large only for small $v$.

The remaining task is to bound $\ell(n,M)$ below, uniformly in $M$. We do this from scratch in
Section~\ref{sec:region}, by exhibiting an explicit interval inside the uncovered region; the following
observation reduces every $k$ to the case of a single deleted speed.

\begin{lemma}[Proved; Monotonicity]\label{lem:mono}
If $M'\subseteq M$ then $U(n,M')\subseteq U(n,M)$; hence for nonempty $M$,
$\ell(n,M)\ge\max_{b\in M}\ell(n,\{b\})$.
\end{lemma}

\begin{theorem}[Proved; Uniform component bound]\label{thm:uniform}
Let $n\ge3$ and let $M\subseteq[n-1]$ be nonempty.
\begin{enumerate}
\item[(i)] $\ell(n,M)\ \ge\ \dfrac{1}{2n^{3}}$, with an elementary and explicit proof.
\item[(ii)] If $M$ contains an element coprime to $n$, or an element $b$ with $2b>n-1$, then
  $\ell(n,M)\ \ge\ \dfrac{1}{n^{2}\,(2\log n+2)}$, again explicitly.
\item[(iii)] $\ell(n,M)\ \ge\ c/(n^{2}\log n)$ for an absolute constant $c$. Here $c$ is not
  explicit: the remaining case rests on Iwaniec's bound for Jacobsthal's function, whose implied
  constant is effective but has not been computed.
\end{enumerate}
\end{theorem}

In the coprime case the truth is sharper, and the sharp constant is attained.

\begin{theorem}[Proved; Coprime component bound]\label{thm:unit}
Let $n\ge2$, let $1\le a<n$ be coprime to $n$, and let $j_+\in[1,n)$ be the unique speed with
$j_+a\equiv 1\pmod n$. If $j_+\in M$ then
\[
  \Bigl[\tfrac{a}{n}-\tfrac{1}{n(n-1)},\ \tfrac{a}{n}\Bigr]\ \subseteq\ U(n,M),
  \qquad\text{so}\qquad \ell(n,M)\ \ge\ \frac{1}{n(n-1)} .
\]
In particular this holds whenever $M$ contains an element coprime to $n$, and for every nonempty $M$
when $n$ is prime. The constant cannot be improved: Proposition~\ref{prop:attained} exhibits
equality, computationally.
\end{theorem}

Combining Theorems~\ref{thm:cover}, \ref{thm:uniform} and \ref{thm:unit}:

\begin{corollary}[Proved]\label{cor:main}
Let $V$ be a $k$-speed modification with $2k<n$ and $\LR(V)\le 1/n$. Then
\[
  \min R\ \le\ \frac{4k\,n^{3}}{n-2k},
\]
and if $M$ contains an element coprime to $n$,
\[
  \min R\ \le\ \frac{2k\,n(n-1)}{n-2k}.
\]
For $n$ prime the second bound is unconditional in $M$. For $k=1$ we have $\min R=\max R$, so these
bound the inserted speed outright.
\end{corollary}

\begin{proof}
Immediate from Theorems~\ref{thm:cover}, \ref{thm:uniform} and \ref{thm:unit}: the first supplies
the covering inequality, the second the component bound valid for every $M$, and the third the
sharper bound when some element of $M$ is coprime to $n$. For $n$ prime every element of $M$ not
divisible by $n$ is coprime to it, and $M\subseteq[n-1]$, so the coprime case always applies.
\end{proof}

\begin{remark}\label{rem:notO}
We state inequalities rather than asymptotics deliberately. Definition~\ref{def:mod} permits
$n-2k=1$, and there the second bound of Corollary~\ref{cor:main} is $\Theta(n^{3})$, not
$\Theta(kn)$, so an assertion of the form $\min R=O(kn)$ is not available from these bounds. We do
not claim it is false: we know of no family realising $\min R\gg kn$, and a search over tight
modifications with $n-2k=1$ found $\max\min R=6$ at $n=5$, $k=2$ and $8$ at $n=7$, $k=3$. That search
is finite for a reason internal to this paper: Theorem~\ref{thm:cover} bounds the smallest inserted
speed, and applying it again to the region left after that speed is removed bounds the next, a branch
terminating as soon as the region has no component of positive length. At such a leaf the branch is
scored without further search, provided the region is nonempty, which holds in every case we
searched and which for fewer than $n-1$ speeds is an instance of the conjecture itself. The region is
then a finite set of rationals; any $t=p/q$ in it, in lowest terms, has $q\ge2$ because
$0\in B_j$ for every present speed $j$ and $[n-1]\setminus M\ne\emptyset$. One may take every remaining speed in the class
$v\equiv c\pmod q$ with $cp\equiv\lfloor q/2\rfloor$, giving $\nrm{vt}=\lfloor q/2\rfloor/q\ge1/3$;
infinitely many such $v$ exist, so $t$ survives and $\min R$ is already determined. Under the additional
hypothesis $2k\le(1-\varepsilon)n$ one does get $\min R=O_\varepsilon(kn)$ in the coprime case,
$O_\varepsilon(kn^{2})$ in general, and $O_\varepsilon(kn\log n)$ from
Theorem~\ref{thm:uniform}(iii).
\end{remark}

Bounding $\ell(n,M)$ from below by exhibiting an interval, as above, leaves open how much of the
uncovered region those intervals miss. They miss nothing: the region has no freedom at all in where
it can sit. The reason is Dirichlet's approximation theorem, which for every real $t$ produces a
denominator $q\le n-1$ approximating $t$ to within $1/(qn)$. If $t$ is uncovered then that
denominator cannot be a present speed, so it was deleted.

\begin{theorem}[Proved; Localization]\label{thm:localize}
Let $n\ge2$ and $M\subseteq[n-1]$, and let $t\in\R$ satisfy $\nrm{jt}>1/n$ for every present $j$.
Then there are $q\in M$ and $a\in\Z$ with $\gcd(a,q)=1$ and
\[
  \Bigl|\,t-\frac aq\,\Bigr|\ \le\ \frac1{qn}.
\]
\end{theorem}

Since there are $\sum_{q\in M}q$ fractions with denominator in $M$, each carrying an interval of
length $2/(qn)$, this bounds the whole region at once, uniformly in how large the deleted speeds
are.

\begin{corollary}[Proved]\label{cor:measure}
$\lambda\bigl(U(n,M)\bigr)\ \le\ 2|M|/n$.
\end{corollary}

More is true. Each component of $U(n,M)$ not only lies near such a fraction $a/q$ but runs back to
it, so the component is cut out by an explicit minimum over the present speeds. That is
Theorem~\ref{thm:exact}: writing $w(n,b,a,M)$ for the window at $a/b$ of
Definition~\ref{def:window}, and provided $M\ne[n-1]$ so that some speed is present,
\[
  \ell(n,M)\ =\ \max_{b\in M}\ \max_{\substack{0\le a<b\\ \gcd(a,b)=1}} w(n,b,a,M),
\]
an identity costing $O\bigl(n\sum_{b\in M}\varphi(b)\bigr)$ operations rather than an
$O(n^{2}\log n)$ sweep. Feeding it into
Theorem~\ref{thm:uncond} gives Corollary~\ref{cor:window}, an unconditional criterion which, unlike
Corollary~\ref{cor:uncond}, loses nothing: over all $k$-subsets at $n=22,\dots,42$ and $k=1,2,3$ it
certifies exactly the subsets for which the criterion of Theorem~\ref{thm:uncond} holds at all.

The coprimality hypothesis in Theorem~\ref{thm:unit} cannot be dropped, and the witness is the
Goddyn--Wong set recorded in version~1 of \cite{zhang} (it is not in version~2): $\{1,\dots,29,31,90\}$ has $n=32$,
$M=\{30\}$, $R=\{90\}$ and $\gcd(30,32)=2$. Here $\ell(32,\{30\})=1/2576$, the bound of
Theorem~\ref{thm:cover} reads $\min R\le 171.73$ and holds, while the coprime bound of
Corollary~\ref{cor:main} reads $\min R\le 66.13$ and fails.

\section{The covering bound}

\begin{lemma}[Proved]\label{lem:strip}
Let $C$ be an interval of length $\ell$, let $n\ge2$ and let $v\ge1$ be an integer. Then
$\lambda(B_v\cap C)\le 2\ell/n+2/(vn)$, where $\lambda$ is Lebesgue measure.
\end{lemma}

\begin{proof}
Work with the $1$-periodic extension of $B_v$ to $\R$, which is legitimate since intersecting back
with $[0,1)$ only decreases measure. It is invariant under translation by $1/v$, and each period
carries measure exactly $2/(vn)$:
the map $t\mapsto vt$ sends a window of length $1/v$ bijectively onto a window of length $1$, and
$\lambda\{u\in[c,c+1):\nrm{u}<1/n\}=2/n$. Put $m=\lceil \ell v\rceil$ and $C=[x,x+\ell]$. Then
$C\subseteq[x,x+m/v]$, a union of $m$ windows of length $1/v$, so
$\lambda(B_v\cap C)\le m\cdot 2/(vn)\le(\ell v+1)\cdot 2/(vn)=2\ell/n+2/(vn)$.
\end{proof}

\begin{proof}[Proof of Theorem~\ref{thm:cover}]
First, $\LR(V)\le 1/n$ forces $S=\{t:\nrm{vt}\ge1/n\text{ for all }v\in V\}$ to be a null set: for
$t\in S$ the minimum $\min_{v\in V}\nrm{vt}$ is at least $1/n$ and, by hypothesis, at most $1/n$, so
it equals $1/n$ and some $v\in V$ has $\nrm{vt}=1/n$. Hence
$S\cap[0,1)\subseteq\bigcup_{v\in V}\{t\in[0,1):\nrm{vt}=1/n\}$, a finite set. So $U(n,M)$ is covered by
$\bigcup_{v\in R}B_v$ up to measure zero.

Let $C\subseteq U(n,M)$ be a component of length $\ell=\ell(n,M)>0$. Then
$\ell\le\sum_{v\in R}\lambda(B_v\cap C)$, so Lemma~\ref{lem:strip} gives
\[
  \ell\;\le\;\sum_{v\in R}\Bigl(\frac{2\ell}{n}+\frac{2}{vn}\Bigr)
   \;=\;\frac{2k\ell}{n}+\frac2n\sum_{v\in R}\frac1v ,
\]
whence $\ell(n-2k)\le 2\sum_{v\in R}1/v$. Since $|R|=k$, $\sum_{v\in R}1/v\le k/\min R$.
\end{proof}

\begin{remark}
The hypothesis $v\ge n$ of Definition~\ref{def:mod} is not used in this proof; Lemma~\ref{lem:strip}
needs only $v\ge1$. It is needed for $|V|=n-1$, not for the argument.
\end{remark}

\section{An unconditional case}\label{sec:uncond}

Theorem~\ref{thm:cover} says that covering forces $\sum_{v\in R}1/v$ to be large. Since every
inserted speed is at least $n$, that sum is at most $k/n$. When the two are incompatible the
covering hypothesis is untenable and no $R$ whatsoever can cover, so $V$ is lonely with no bound on
any inserted speed and no appeal to anything further.

\begin{theorem}[Proved; Unconditional loneliness]\label{thm:uncond}
Let $n\ge3$ and $k\ge1$ with $2k<n$, and let $M\subseteq[n-1]$ with $|M|=k$. If
\[
  \ell(n,M)\ >\ \frac{2k}{n\,(n-2k)},
\]
then $\LR\bigl(([n-1]\setminus M)\cup R\bigr)\ge 1/n$ for \emph{every} admissible $R$, that is, for
every $R$ as in Definition~\ref{def:mod}.
\end{theorem}

\begin{proof}
If $\LR(V)<1/n$ then Theorem~\ref{thm:cover} gives
$\sum_{v\in R}1/v\ge\ell(n,M)(n-2k)/2>k/n$. But $|R|=k$ and every $v\in R$ satisfies $v\ge n$, so
$\sum_{v\in R}1/v\le k/n$, a contradiction, since Theorem~\ref{thm:cover} applies as soon as $\LR(V)\le1/n$. Hence
$\LR(V)\ge1/n$.
\end{proof}

Combined with the lower bounds of Section~\ref{sec:region} this settles an explicit subclass
outright. Feeding Proposition~\ref{prop:missing}'s bound $\ell(n,\{b\})\ge(2s-j_0)/(2nbj_0)$, valid
for $2b\le n-1$, and Proposition~\ref{prop:twosided}'s bound
$\ell(n,\{b\})\ge\frac{s}{nb}(\frac1{j_0}+\frac1{n-1})$, valid for $2b\ge n$, into the criterion, and
using Lemma~\ref{lem:mono}:

\begin{corollary}[Proved]\label{cor:uncond}
Let $n\ge3$ and $k\ge1$ with $2k<n$, and let $M\subseteq[n-1]$ with $|M|=k$. Write $s=n-b$, and for $b\ge2$ let $j_0=j_0(n,b)$ be as in Section~\ref{sec:region} (case (a)
below does not use it). Suppose $M$ contains a speed $b$ with either
\begin{enumerate}
\item[(a)] $b=1$ and $n>6k$; or
\item[(b)] $2\le b$, $2b\le n-1$ and $(2s-j_0)(n-2k)>4kb\,j_0$; or
\item[(c)] $2\le b$, $2b>n-1$ and $s(n-2k)(j_0+n-1)>2kb\,j_0(n-1)$.
\end{enumerate}
Then $\LR(([n-1]\setminus M)\cup R)\ge1/n$ for every admissible $R$.
\end{corollary}

\begin{proof}
Case (a) is separate because Proposition~\ref{prop:missing} requires $b\ge2$: the integer $a$ there
lies in $[1,b)$, which is empty at $b=1$. It is handled directly, and the argument is that of
\cite[proof of Thm.~2.3, case $r=1$]{gw}. Every present speed $j$ satisfies
$2\le j\le n-1$, so for $t\in[\tfrac1{2n},\tfrac1n]$ we have
$jt\in[\tfrac{j}{2n},\tfrac{j}{n}]\subseteq[\tfrac1n,1-\tfrac1n]$, whence $\nrm{jt}\ge1/n$. Thus
$[\tfrac1{2n},\tfrac1n]\subseteq U(n,\{1\})$ and $\ell(n,\{1\})\ge\tfrac1{2n}$, which exceeds
$2k/(n(n-2k))$ exactly when $n>6k$. (In fact $\ell(n,\{1\})=1/(2n)$; we verified this for
$3\le n\le44$.)

In case (b), Proposition~\ref{prop:missing} gives $\ell(n,\{b\})\ge(2s-j_0)/(2nbj_0)$; in case (c),
Proposition~\ref{prop:twosided} gives $\ell(n,\{b\})\ge\frac{s}{nb}(\frac1{j_0}+\frac1{n-1})$. The
stated inequalities say exactly that these exceed $2k/(n(n-2k))$. Lemma~\ref{lem:mono} raises the bound from $\ell(n,\{b\})$ to $\ell(n,M)$, and
Theorem~\ref{thm:uncond} applies.
\end{proof}

\begin{remark}\label{rem:closed}
Condition (b) admits a closed form at the cost of some strength. Since $j_0\le n-1$ one has
$(2s-j_0)/j_0\ge(n-2b+1)/(n-1)$, so the inequality clause of (b) follows from
$4kb(n-1)<(n-2b+1)(n-2k)$, and that in turn from $8kb\le n$ together with $n>6k-2$. The side clauses
of (b) come along too when $b\ge2$: $8kb\le n$ gives $2b\le n/(4k)\le n/4\le n-1$. (At $b=1$ the
hypotheses land in case (a) instead.) The derivation: the first gives $4kb(n-1)\le n(n-1)/2$, while
$2b\le n/(4k)\le n/4$ gives $(n-2b+1)(n-2k)\ge\tfrac34n(n-2k)$, and $n(n-1)/2<\tfrac34n(n-2k)$
reduces to $n>6k-2$. We verified $4kb(n-1)<(n-2b+1)(n-2k)$ directly on all $4\,338\,029$ triples
with $b\ge2$, $n\le4000$ and $k\le60$ admitted by those hypotheses, with no exception. (Admitting
$b=1$, which the previous sentence sends to case (a) instead, the count is $4\,563\,449$ and there
is still no exception.)

The step $j_0\le n-1$ is wasteful for mid-range $b$, where $j_0$ exceeds $s$ only by the gap to the
next integer coprime to $b$. If that gap is $o(s)$ then (b) reads $n-2k>4kb$ asymptotically, i.e.
$b<(n-2k)/(4k)$, roughly twice the range of the closed form. We do not use this, since bounding the
gap needs the input of Theorem~\ref{thm:uniform}(iii); conditions (b) and (c) are exact and
checkable as they stand.
\end{remark}

How much this settles, counted exactly over \emph{all} $k$-subsets $M$. Over $22\le n\le38$ the
criterion of Theorem~\ref{thm:uncond} holds for $36.00\%$ to $44.44\%$ of subsets when $k=1$, for
$17.66\%$ to $25.08\%$ when $k=2$, and for $11.88\%$ to $19.56\%$ when $k=3$. Its reach falls away
for small $n$: at $n=10$ it settles $22.22\%$ of subsets when $k=1$ and nothing at all when $k=2$ or
$k=3$, the first $n$ at which it is satisfiable being $7,11,15$ for $k=1,2,3$.
Corollary~\ref{cor:uncond} recovers most of that. At $n=22,26,30,34,38,42$, with the criterion's own
reach printed alongside so that the comparison is on the same $n$:
\[
\begin{array}{lcccccc}
\hline
n & 22 & 26 & 30 & 34 & 38 & 42\\
\hline
\text{criterion, }k=1 & 38.10 & 36.00 & 41.38 & 42.42 & 37.84 & 41.46\\
\text{Cor.~\ref{cor:uncond}, }k=1 & 33.33 & 32.00 & 37.93 & 33.33 & 35.14 & 36.59\\
\text{ratio} & 0.875 & 0.889 & 0.917 & 0.786 & 0.929 & 0.882\\[0.6ex]
\text{criterion, }k=2 & 21.43 & 19.33 & 23.89 & 21.78 & 24.77 & 22.44\\
\text{Cor.~\ref{cor:uncond}, }k=2 & 18.57 & 15.67 & 19.95 & 17.61 & 20.72 & 18.78\\
\text{ratio} & 0.867 & 0.811 & 0.835 & 0.808 & 0.836 & 0.837\\[0.6ex]
\text{criterion, }k=3 & 14.44 & 13.13 & 11.88 & 18.37 & 16.82 & 15.55\\
\text{Cor.~\ref{cor:uncond}, }k=3 & 14.29 & 12.00 & 10.34 & 17.61 & 15.77 & 14.27\\
\text{ratio} & 0.990 & 0.914 & 0.870 & 0.959 & 0.938 & 0.918\\
\hline
\end{array}
\]
The closed form of Remark~\ref{rem:closed} is weaker, certifying
$9.52,12.00,10.34,12.12,10.81,12.20$ per cent at $k=1$ over the same $n$. The residual gap is closed
outright in Section~\ref{sec:localize}: Corollary~\ref{cor:window} replaces the case analysis by a
single explicit window and matches the criterion in every cell of the table above.

The two-sided window of Proposition~\ref{prop:twosided} is what makes case (c) useful at $k=1$. With
the one-sided bound of Proposition~\ref{prop:missing} in its place, case (c) is provably empty at
every $k$: in that regime $2b\ge n$ and $j_0\ge s$ give $2kbj_0\ge kns\ge ns>(n-2k)s$, so the
one-sided inequality $s(n-2k)>2kbj_0$ has no solutions at all, and the $k=1$ reach falls to
$19.05,20.00,20.69,21.21,21.62,21.95$ per cent. Even with the two-sided window, case (c) remains
empty for every $k\ge2$; see Remark~\ref{rem:cempty}.

The reason the criterion fires so often is that $\ell(n,\{b\})$ is of order $1/(2bn)$ for small $b$,
not of order $1/n^{2}$: at $n=40$ the values of $\ell\,n^{2}$ for $b=1,\dots,9$ are
$20.000,9.487,6.667,4.730,3.778,3.143,2.857,2.348,2.222$, in close agreement with $n/(2b)$. This is
what Proposition~\ref{prop:missing} predicts, since $s\approx n$ and $j_0\approx n$ there.

\section{The uncovered region for general \texorpdfstring{$M$}{M}}\label{sec:region}

\begin{proof}[Proof of Lemma~\ref{lem:mono}]
$[n-1]\setminus M\subseteq[n-1]\setminus M'$, so the constraints defining $U(n,M)$ are a subfamily of
those defining $U(n,M')$. Taking $M'=\{b\}$ for $b\in M$ gives the second claim.
\end{proof}

By Lemma~\ref{lem:mono} it suffices to exhibit a long interval inside $U(n,\{b\})$ for a single $b$.
We do this directly. Fix $2\le b<n$, put $s=n-b\ge1$ and
\[
  j_0\;=\;j_0(n,b)\;=\;\min\{\,j\ge s:\ \gcd(j,b)=1\,\}.
\]
The window $[s,s+b-1]$ consists of $b$ consecutive integers, one of them congruent to $1$ modulo $b$,
so $j_0$ exists and $s\le j_0\le s+b-1=n-1$. In particular $j_0$ is a speed of $[n-1]$, and
$j_0\ne b$ because $\gcd(j_0,b)=1$ and $b\ge2$.

\begin{proposition}[Proved; Missing-speed interval]\label{prop:missing}
Let $2\le b<n$ and let $a$ be the unique integer with $1\le a<b$ and $j_0\,a\equiv-1\pmod b$; it
exists because $\gcd(j_0,b)=1$, and $a\ne0$ because $\gcd(a,b)=1$. Put
\[
  \delta_0=\begin{cases}\dfrac{1}{2bn}, & 2b\le n-1,\\[2mm] 0, & 2b>n-1.\end{cases}
\]
Then the interval below is nonempty: if $2b\le n-1$ then $s\ge b+1$, so $j_0\le s+b-1\le 2s-2$ and
$\delta_0<s/(nbj_0)$; while if $2b>n-1$ then $\delta_0=0$. (The claim $j_0\le 2s$ is false in the
second regime: at $n=32$, $b=30$ one has $s=2$ and $j_0=7$.) Moreover
$\bigl[\tfrac ab+\delta_0,\ \tfrac ab+\tfrac{s}{nbj_0}\bigr]\subseteq U(n,\{b\})\subseteq[0,1)$; hence
\[
  \ell(n,\{b\})\ \ge\ \frac{s}{n\,b\,j_0}-\delta_0
  \;=\;\begin{cases}\dfrac{2s-j_0}{2\,n\,b\,j_0}, & 2b\le n-1,\\[2mm]
                    \dfrac{s}{n\,b\,j_0}, & 2b>n-1.\end{cases}
\]
\end{proposition}

\begin{proof}
Let $\delta\in[\delta_0,\,s/(nbj_0)]$ and $t=a/b+\delta$. Note $\delta\le 1/(nb)$, since $j_0\ge s$.
Let $j\in[1,n)$ with $j\ne b$.

First, $t\in[0,1)$: since $1\le a<b$ we have $t\ge a/b\ge 1/b>0$, and
$t\le\frac{b-1}{b}+\frac{1}{nb}=1-\frac{n-1}{nb}<1$.

\emph{Multiples of $b$.} If $b\mid j$, write $j=mb$; then $m\ne0$ and $m\ne1$, so $m\ge2$ and
$j\ge 2b$, which forces $2b\le n-1$ and hence $\delta_0=1/(2bn)$. Now $j\cdot(a/b)=ma\in\Z$, so
$\nrm{jt}=\nrm{j\delta}$, and
\[
  j\delta\ \ge\ 2b\cdot\frac{1}{2bn}=\frac1n,
  \qquad
  j\delta\ \le\ (n-1)\cdot\frac{1}{nb}\ \le\ \frac{n-1}{n}=1-\frac1n .
\]
A real number in $[1/n,1-1/n]$ is at distance at least $1/n$ from every integer.

\emph{Non-multiples.} If $b\nmid j$, put $r=ja\bmod b$; then $r\ne0$ because $\gcd(a,b)=1$, so
$1\le r\le b-1$. Writing $ja=bQ+r$ gives $jt=Q+\bigl(\tfrac rb+j\delta\bigr)$, and
$\tfrac rb+j\delta\ge\tfrac1b\ge\tfrac1n$. For the upper edge we must check
$j\delta\le\frac{b-r}{b}-\frac1n$, and it suffices to check it at $\delta=s/(nbj_0)$.
If $r=b-1$ then $j\equiv -a^{-1}\equiv j_0\pmod b$; the elements of that class below $n$ are
$j_0,\ j_0-b,\ j_0-2b,\dots$, since $j_0+b\ge s+b=n$, so $j\le j_0$ and
\[
  j\delta\ \le\ j_0\cdot\frac{s}{nbj_0}=\frac{s}{nb}=\frac1b-\frac1n=\frac{b-r}{b}-\frac1n .
\]
If $r\le b-2$ then, using $j\le n-1$ and $s\le j_0$,
\[
  j\delta\ \le\ \frac{(n-1)s}{nbj_0}\ \le\ \frac{n-1}{nb}
  \ \le\ \frac{2n-b}{nb}\ =\ \frac2b-\frac1n\ \le\ \frac{b-r}{b}-\frac1n ,
\]
the middle inequality because $b\le n+1$. So no residue class with $b-r\ge2$ is ever binding.
\end{proof}

\begin{proposition}[Proved; Two-sided window]\label{prop:twosided}
Let $2\le b<n$ with $2b>n-1$, and let $a$, $j_0$ be as in Proposition~\ref{prop:missing}. Then no
multiple of $b$ other than $b$ itself is a speed of $[n-1]$, the window extends on both sides of
$a/b$, and
\[
  \ell(n,\{b\})\ \ge\ \frac{s}{n\,b}\Bigl(\frac1{j_0}+\frac1{n-1}\Bigr).
\]
\end{proposition}

\begin{proof}
Since $2b>n-1$, any multiple $mb$ with $m\ge2$ satisfies $mb\ge 2b\ge n$, so the only present speeds
are non-multiples of $b$, and Proposition~\ref{prop:missing} applies with $\delta_0=0$ on the right.
On the left, take $t=a/b-\delta$ with $\delta\ge0$. For a present $j$ with $r=ja\bmod b$, the
argument of Proposition~\ref{prop:missing} applies with the roles of $r$ and $b-r$ exchanged: the
binding class is now $r=1$, that is $j\equiv a^{-1}\pmod b$, and the binding member is the largest
such $j$ below $n$, which is at most $n-1$. Hence $\delta\le s/(nb(n-1))$ is admissible, and the
interval $[\,a/b-s/(nb(n-1)),\ a/b+s/(nbj_0)\,]$ lies in $U(n,\{b\})$.
\end{proof}

We verified this bound against the exhaustively computed $\ell(n,\{b\})$ for all $9999$ pairs with
$2\le b<n\le200$ and $2b\ge n$: no violation. Equality holds throughout the family $b=n/2$, $n$ even, which we verified for all even $6\le n\le250$; over all admissible pairs with $n\le N$ the proportion of exact cases falls from $14.3\%$ at $N=30$ to $1.8\%$ at $N=250$. The bound is exact, for instance
$\ell(8,\{4\})=3/70=1/40+1/56$ and $\ell(32,\{16\})=1/544+1/992$.

\begin{remark}\label{rem:zhang4}
Proposition~\ref{prop:twosided} lives in the regime $2b\ge n$, and there it is
\cite[Thm.~1.4]{zhang} weakened by one substitution. Zhang computes
$\ell(n,\{b\})=\frac{s}{bn}\max_u\bigl(\frac1{F(u)}+\frac1{F(b-u)}\bigr)$ with $F(x)$ the largest
$i\le n-1$ congruent to $x$ modulo $b$; for $x$ coprime to $b$ that representative is the unique one
in $[s,n-1]$, so $\min_u F(u)$ is exactly our $j_0$, and his $I(n,r)$ is our $j_0$. Replacing
$F(b-u)$ by the cruder $n-1$ gives Proposition~\ref{prop:twosided}, which is why ours is an
inequality where his is an identity. A reader wanting the sharp statement in this regime should take
it from \cite{zhang}; we keep the weaker form only because it is what the Lean development proves.
\end{remark}

\begin{remark}\label{rem:zhang}
Proposition~\ref{prop:missing} is a lower bound only, which is all that is needed below. Zhang
\cite[Thm.~1.4]{zhang} computes $\ell(n,\{b\})$ exactly, and in the regime $2b\le n-1$ his value
agrees with ours, so the bound is sharp there.
\end{remark}

\begin{proof}[Proof of Theorem~\ref{thm:uniform}]
By Lemma~\ref{lem:mono} take $M=\{b\}$. If $b=1$ then $\gcd(b,n)=1$ and $j_+=1=b\in M$ for $a=1$,
so Theorem~\ref{thm:unit} gives $\ell\ge 1/(n(n-1))$, which is stronger than every claim below.
Assume $b\ge2$ and apply Proposition~\ref{prop:missing}. The two regimes are separated by $2b$
against $n-1$, not by $b$ against $n/2$.

\emph{(i).} If $2b>n-1$ then $\ell\ge s/(nbj_0)\ge 1/(n\cdot n\cdot n)=1/n^{3}$, using $s\ge1$,
$b<n$ and $j_0\le n-1$. If $2b\le n-1$ then $s=n-b\ge b+1$, so $j_0\le s+b-1\le 2s-2$, giving
$2s-j_0\ge2$ and
\[
  \ell\ \ge\ \frac{2s-j_0}{2nbj_0}\ \ge\ \frac{1}{nbj_0}\ \ge\ \frac{2}{n(n-1)^{2}}\ \ge\ \frac{1}{2n^{3}},
\]
using $b\le(n-1)/2$ and $j_0\le n-1$, and $4n^{2}\ge(n-1)^{2}$ for the last step.

\emph{(ii).} Suppose $2b>n-1$, so $\ell\ge s/(nbj_0)$. Let $P=P(b)$ and let $P^m$ be the least power
of $P$, with $m\ge0$ allowed, satisfying $P^m\ge s$. Then $P^m<Ps$ (for $s=1$ take $m=0$) and
$\gcd(P^m,b)=1$, so $j_0\le P^m<Ps$ and $\ell> s/(nb\cdot Ps)=1/(nbP)>1/(n^{2}P)$. It remains to
check $P<2\log b+2$. If $P\ge13$ then every prime at most $11$ divides $b$, so
$b\ge e^{\theta(P-1)}>e^{(P-1)/2}$, where $\theta(x)=\sum_{p\le x}\log p$ and $\theta(x)>x/2$ for
$x\ge11$ \cite{rs};
hence $P<2\log b+1$. For $P\le11$ the least admissible $b$ is $3,2,6,30,210$ for $P=2,3,5,7,11$, and
since $2\log b+2$ increases with $b$ it suffices to check those: the values are
$4.20,3.39,5.58,8.80,12.69$, each larger than $P$. If instead $M$ contains an element coprime to
$n$, Theorem~\ref{thm:unit} gives $\ell\ge 1/(n(n-1))\ge 1/(n^{2}(2\log n+2))$.

\emph{(iii).} By (ii) the cases $2b>n-1$ and $\gcd(b,n)=1$ are done, so there remains $2b\le n-1$
with $\gcd(b,n)>1$, where $s\ge b+1$. Let $g(b)$ be Jacobsthal's function, the largest gap between
consecutive integers coprime to $b$; then $j_0\le s+g(b)-1$. By Iwaniec \cite{iwaniec},
$g(b)\ll\log^{2}b$, so there is $N_0$ with $g(b)\le s/2$ for all $n\ge N_0$ in this range. For such
$n$, $2s-j_0\ge s/2$ and $j_0\le 3s/2$, so
$\ell\ge \frac{s/2}{2nb\cdot 3s/2}=\frac{1}{6nb}\ge\frac{1}{3n^{2}}$. For $3\le n<N_0$, part (i)
gives $\ell\ge 1/(2n^{3})$, hence $\ell\,n^{2}\log n\ge \log 3/(2N_0)$. Taking $c$ to be the least of
$1/3$, $\log 3/(2\log 3+2)$ and $\log 3/(2N_0)$ proves (iii); the second of these is the constant
obtained from (ii), since converting $1/(n^{2}(2\log n+2))$ to the shape $c/(n^{2}\log n)$ requires
$c\le\log n/(2\log n+2)$, whose infimum over $n\ge3$ is $\log 3/(2\log 3+2)=0.2617\ldots$. The
dependence on $N_0$, hence on Iwaniec's implied constant, is why $c$ is not explicit.
\end{proof}

\section{Localization of the uncovered region}\label{sec:localize}

This section is expository. For $M=\{r\}$ the localization below is in \cite[\S2: Lemmas~2.1 and~2.2 and the first paragraph of the proof of Lemma~2.4; the case $2r>n-1$ is Lemma~2.5]{zhang}, version~2:
Dirichlet's theorem with $N=n-1$ forces $\nrm{rt}\le1/n$ on $U(n,r)$, whence
$U(n,r)\subseteq\bigcup_p[p/r-1/(rn),\,p/r+1/(rn)]$, coprimality reduction included. Replacing
$\{r\}$ by $M$ changes nothing in that argument. Bohman and Peng \cite[\S3]{bp} use the same idea
constructively, taking a witness time whose denominator is built from deleted speeds. We note also
that Dirichlet is not needed for the containment: since $\LR([n-1])=1/n$, the set on which every
$j\in[n-1]$ has $\nrm{jt}\ge1/n$ is the finite set $\{a/n:\gcd(a,n)=1\}$ (see
the opening line of the proof of \cite[Prop.~10.4]{kravitz}), so off a finite set $U(n,M)\subseteq\bigcup_{b\in M}B_b$ in one line,
which also gives Corollary~\ref{cor:measure} directly. The measure of the uncovered region as an
object of study is \cite{cusick74}, and measure together with component count already appears in
\cite[Lemma~1]{cp}. We record all of it here because the formalization and Section~\ref{sec:comp}
consume it.

Sections~\ref{sec:uncond} and~\ref{sec:region} bound $\ell(n,M)$ from below by exhibiting an
interval inside $U(n,M)$. This section determines $\ell(n,M)$ exactly. The mechanism is Dirichlet's
approximation theorem: it forces every uncovered time to sit close to a fraction whose denominator
was deleted, and there is nowhere else for the uncovered region to be.



\begin{proof}[Proof of Theorem~\ref{thm:localize}]
Dirichlet's approximation theorem with $N=n-1$ gives an integer $q$ with $1\le q\le n-1$ and
$\nrm{qt}\le 1/(N+1)=1/n$. If $q$ were present the hypothesis would force $\nrm{qt}>1/n$, so
$q\in M$. Taking $a$ nearest to $qt$ turns $\nrm{qt}\le1/n$ into $|t-a/q|\le1/(qn)$.

For coprimality, write $d=\gcd(a,q)$, $a=da'$, $q=dq'$. Then $|t-a'/q'|=|t-a/q|\le 1/(qn)\le
1/(q'n)$, so $\nrm{q't}\le 1/n$ with $1\le q'\le n-1$, and the same argument puts $q'$ in $M$.
\end{proof}

The first consequence bounds the size of the uncovered region, uniformly in which speeds were
deleted and how large they are.



\begin{proof}[Proof of Corollary~\ref{cor:measure}]
The complement of $U(n,M)$ in $[0,1)$ is a finite union of open intervals, so $U(n,M)$ differs from
its interior by a finite set. Every interior point satisfies $\nrm{jt}>1/n$ for all present $j$:
if some present $j$ had $\nrm{jt}=1/n$ then $t$ would lie on the boundary of $B_j$ and hence not be
interior to the complement of $B_j$. Theorem~\ref{thm:localize} therefore covers the interior by the
intervals $[a/q-1/(qn),\,a/q+1/(qn)]$ with $q\in M$ and $0\le a<q$. There are at most
$\sum_{q\in M}q$ of them, of length $2/(qn)$ each, for a total of
$\sum_{q\in M}q\cdot 2/(qn)=2|M|/n$.
\end{proof}

The bound is of the right order: over all $M$ with $|M|=k$ for $k=1,2,3$ and
$n\in\{16,22,28,34,40\}$ the largest value of $\lambda(U(n,M))\big/(2k/n)$ lies between $0.5139$ and
$0.6492$. The maximiser is unique in each of the fifteen cells and is usually a set of primes
just above $n/2$, though not always: at $n=16$, $k=2$ it is $\{9,11\}$, at $n=16$, $k=3$ it is
$\{9,11,13\}$, at $n=28$, $k=3$ it is
$\{1,17,19\}$, and at $n=40$, $k=3$ it is $\{23,25,29\}$.

The second consequence is an exact formula. The point of Theorem~\ref{thm:localize} is not only that
the uncovered region sits near the deleted denominators, but that each of its components runs all
the way back to one of them, so that the component is determined by an explicit minimum.

\begin{definition}\label{def:window}
Let $b\in M$ and let $a$ satisfy $0\le a<b$ and $\gcd(a,b)=1$. For present $j$ write
$r_j=ja\bmod b$, and set
\[
  \delta^{+}=\min_{j\ \text{present}}
    \begin{cases}\dfrac{(b-r_j)/b-1/n}{j}, & b\nmid j,\\[2.2ex]
                 \dfrac{1-1/n}{j}, & b\mid j,\end{cases}
  \qquad
  \delta^{-}=\min_{j\ \text{present}}
    \begin{cases}\dfrac{r_j/b-1/n}{j}, & b\nmid j,\\[2.2ex]
                 \dfrac{1-1/n}{j}, & b\mid j,\end{cases}
\]
\[
  \delta^{0}=\max\Bigl\{\tfrac1{jn}\ :\ j\ \text{present},\ b\mid j\Bigr\},
\]
with $\delta^{0}=0$ when $b$ has no present multiple. The \emph{window} at $a/b$ is
\[
  w(n,b,a,M)=\begin{cases}
    \delta^{+}+\delta^{-}, & \delta^{0}=0,\\[0.6ex]
    \max\{\delta^{+},\delta^{-}\}-\delta^{0}, & \delta^{0}>0,
  \end{cases}
\]
read as $0$ if the right-hand side is negative.
\end{definition}

\begin{theorem}[Proved; Exact formula]\label{thm:exact}
For every $n\ge3$ and every $M$ with $\emptyset\ne M\subsetneq[n-1]$,
\[
  \ell(n,M)\ =\ \max_{b\in M}\ \max_{\substack{0\le a<b\\ \gcd(a,b)=1}}\ w(n,b,a,M).
\]
The hypothesis $M\ne[n-1]$ is needed: at $M=[n-1]$ there is no present speed, $U(n,M)=[0,1)$ and
$\ell=1$, while $\delta^{+}$ and $\delta^{-}$ are minima over an empty family. Every $M$ arising
from a $k$-speed modification satisfies $2k<n$ and so is covered.
\end{theorem}

\begin{proof}
\emph{The maximum is attained inside $U(n,M)$.} Fix $b$ and $a$ and let $t=a/b+\delta$ with
$\delta^{0}\le\delta\le\delta^{+}$. If $b\mid j$ then $j\,a/b\in\Z$, so $\nrm{jt}=\nrm{j\delta}$;
here $j\delta\ge j\delta^{0}\ge 1/n$ by the definition of $\delta^{0}$, and $j\delta\le
j\delta^{+}\le 1-1/n$ by the second branch of $\delta^{+}$. If $b\nmid j$ then $jt\equiv
r_j/b+j\delta$, which is at least $r_j/b\ge1/b>1/n$ above the integer below it, and at most
$1-1/n$ by the first branch. So $[a/b+\delta^{0},\,a/b+\delta^{+}]\subseteq U(n,M)$, and symmetrically
on the left; when $\delta^{0}=0$ the two sides join at $a/b$.

\emph{Every component of positive length is such a window.} Let $C$ be a component of $U(n,M)$
with $|C|>0$, which exists by Theorem~\ref{thm:uniform}(i), and let $t$ be an interior
point of $C$, so $\nrm{jt}>1/n$ for every present $j$. Theorem~\ref{thm:localize} gives $b\in M$ and
$a$ coprime to $b$ with $|t-a/b|\le 1/(bn)$; assume $t\ge a/b$, the other case being the mirror
image, and put $D=t-a/b$. Since $j\le n-1$ we have $jD\le (n-1)/(bn)<1/b$ for every present $j$.

First, $\delta^{0}\le D$. Indeed if $b\mid j$ for a present $j$ then $\nrm{jt}=\nrm{jD}$, and
$jD<1/b\le1$, so $\nrm{jt}>1/n$ forces $jD>1/n$, that is $D>1/(jn)$; taking the maximum over the
present multiples gives $D\ge\delta^{0}$.

Now let $\delta^{0}\le d\le D$ and $s=a/b+d$. If $b\nmid j$, write $ja=Qb+r_j$ with $0<r_j<b$; then
$js=Q+r_j/b+jd$. Its distance down to $Q$ is $r_j/b+jd\ge r_j/b\ge1/b>1/n$, and its distance up to
$Q+1$ is $1-r_j/b-jd$, which decreases in $d$ and so is at least $1-r_j/b-jD$. Now
$r_j/b+jD<(b-1)/b+1/b=1$ rules out a wrap, so $1-r_j/b-jD$ is one of the two distances from $jt$ to
$\Z$ and hence at least $\nrm{jt}$, which is at least $1/n$ because $t\in U(n,M)$. If instead
$b\mid j$ then $\nrm{js}=\nrm{jd}$, with $jd\ge j\delta^{0}\ge1/n$ and
$jd\le jD\le (n-1)/(bn)\le 1-1/n$.

So no present speed blocks between $a/b+\delta^{0}$ and $t$, whence $[a/b+\delta^{0},t]\subseteq C$.
It remains to identify both endpoints of $C$. On the right, the containment step already gives
$[a/b+\delta^{0},a/b+\delta^{+}]\subseteq U(n,M)$, and at $\delta^{+}$ the minimum defining it is
attained by some present $j$, for which $\nrm{j(a/b+\delta)}<1/n$ for every $\delta$ slightly larger;
so $C$ ends at $a/b+\delta^{+}$. On the left, if $b$ has a present multiple $j$ then
$\nrm{j(a/b+\delta)}=\nrm{j\delta}<1/n$ for $0\le\delta<1/(jn)$, so $C$ starts at $a/b+\delta^{0}$
and $|C|=\delta^{+}-\delta^{0}$; the mirror argument bounds the component on the left of $a/b$ by
$\delta^{-}-\delta^{0}$, and $w$ takes the larger of the two. If $b$ has no present multiple then
$\delta^{0}=0$, the two sides join at $a/b$, and the same argument at $-\delta^{-}$ gives
$C=[a/b-\delta^{-},a/b+\delta^{+}]$, of length $\delta^{+}+\delta^{-}$. Either way
$|C|=w(n,b,a,M)$, and taking $C$ of maximal length gives the reverse inequality.
\end{proof}

\begin{remark}\label{rem:measeq}
Since the components are disjoint, $\lambda(U(n,M))$ is the sum of their lengths. It is
\emph{not} the case that every component is a window: $U(n,M)$ also has degenerate components,
isolated points of the equality locus where some present $j$ attains $\nrm{jt}=1/n$ exactly. At
$n=4$, $M=\{2\}$ there are three components and one window. Over all $65\,504$ nonempty
$M\subsetneq[n-1]$ with $3\le n\le16$, degenerate non-window components occur in $43\,849$ cases.
They carry no measure, so what is true, and what Theorem~\ref{thm:exact} gives, is that every
\emph{positive-length} component is a window; that holds with no exception over the same $65\,504$
cases. Hence $\lambda(U(n,M))$ is the sum of the windows \emph{distinct as intervals}: the
deduplication is load-bearing, and summing over base points without it is wrong in $62\,816$ of
those $65\,504$ cases, because a base point whose $b$ has a present multiple contributes the two
windows $[a/b+\delta^{0},a/b+\delta^{+}]$ and $[a/b-\delta^{-},a/b-\delta^{0}]$ while a component
may abut two different base points. This is bookkeeping, not a theorem, and the pairing of measure
with component count is already in \cite[Lemma~1]{cp}. Verified on all $7\,315$ pairs with
$|M|\le3$, $n\in\{12,16,20,24,28\}$.
\end{remark}

Three points about the formula. It costs $O\bigl(n\sum_{b\in M}\varphi(b)\bigr)$ arithmetic
operations, against $O(n^{2}\log n)$ for sweeping the endpoints $(in\pm1)/(jn)$ directly; for a
single deleted speed $b$ of order $n$ that is a saving of the logarithmic factor only, and the
saving is real only when the deleted speeds are small. It excludes \emph{all} of $M$ from the minima,
which is what Lemma~\ref{lem:mono} cannot do: monotonicity replaces $M$ by a single $\{b\}$ and so
still sees the other deleted speeds as present. And it turns Theorem~\ref{thm:uncond} into a finite,
exactly checkable condition.

\begin{corollary}[Proved; Unconditional criterion, explicit form]\label{cor:window}
Let $n\ge3$ and $k\ge1$ with $2k<n$, and let $M\subseteq[n-1]$ with $|M|=k$. If some $b\in M$ and
some $a$ with $0\le a<b$ and $\gcd(a,b)=1$ satisfy
\[
  w(n,b,a,M)\ >\ \frac{2k}{n(n-2k)},
\]
then $\LR(([n-1]\setminus M)\cup R)\ge 1/n$ for every admissible $R$.
\end{corollary}

\begin{proof}
Immediate from Theorem~\ref{thm:exact} and Theorem~\ref{thm:uncond}.
\end{proof}

Corollary~\ref{cor:window} is stronger than Corollary~\ref{cor:uncond}, and by
Theorem~\ref{thm:exact} it is equivalent to the criterion of Theorem~\ref{thm:uncond} itself, so it
exhausts what the covering bound can reach. The table below is therefore a consistency check on the
implementation rather than independent evidence; we record it because it is what the code computes.
Counted over \emph{all} $k$-subsets $M$, at $n=22,26,30,34,38,42$:
\[
\begin{array}{lcccccc}
\hline
n & 22 & 26 & 30 & 34 & 38 & 42\\
\hline
k=1 & 38.10 & 36.00 & 41.38 & 42.42 & 37.84 & 41.46\\
k=2 & 21.43 & 19.33 & 23.89 & 21.78 & 24.77 & 22.44\\
k=3 & 14.44 & 13.13 & 11.88 & 18.37 & 16.82 & 15.55\\
\hline
\end{array}
\]
The same agreement holds at $n=20,24,28,32,36,40$ for $k=1,2,3$, and the identity
$\ell(n,M)=\max_{b,a}w(n,b,a,M)$ of Theorem~\ref{thm:exact} was checked directly, in exact rational
arithmetic, on all $134\,721$ pairs $(n,M)$ with $|M|\le4$ and $n\in\{16,20,24,28,32,36\}$.

\begin{remark}\label{rem:cempty}
Condition (c) of Corollary~\ref{cor:uncond} is empty for every $k\ge2$. In that regime $2b\ge n$,
and $s\le j_0\le n-1$, so
\[
  s(n-2k)(j_0+n-1)\ \le\ j_0(n-4)\cdot 2(n-1)\ <\ 2n\,j_0(n-1)\ \le\ 2kb\,j_0(n-1),
\]
the last step using $k\ge2$ and $2b\ge n$. Condition (c) therefore contributes only at $k=1$, where
it fires $331\,291$ times for $n\le2000$; for $2\le k\le10$ and $n\le2000$ it fires not at all. This
is why the $k=2$ and $k=3$ figures for conditions (a)--(c) coincide with those for (a)--(b) alone.
Corollary~\ref{cor:window} is not subject to this: it fires at every $k$ for which the criterion
does.
\end{remark}

\section{The coprime case}

\begin{proof}[Proof of Theorem~\ref{thm:unit}]
Since $1\le a<n$ we have $a/n-\frac1{n(n-1)}\ge0$ and $a/n\le\frac{n-1}{n}<1$, so the interval lies
in $[0,1)$. Let $j\in[1,n)$ with $j\ne j_+$, let $\varepsilon\in[0,\frac{1}{n(n-1)}]$ and put
$t=a/n-\varepsilon$. Write $r=ja\bmod n$. As $\gcd(a,n)=1$ and $1\le j<n$ we have $n\nmid ja$, so
$r\ne0$; and $r=1$ would give $j\equiv a^{-1}$, that is $j=j_+$, so $r\ne1$. Hence $2\le r\le n-1$.
Also $j\varepsilon\le(n-1)\cdot\frac1{n(n-1)}=\frac1n$. Writing $ja=nq+r$,
\[
  jt=q+\Bigl(\frac rn-j\varepsilon\Bigr),\qquad
  \frac rn-j\varepsilon\in\Bigl[\frac{r-1}{n},\frac rn\Bigr]\subseteq\Bigl[\frac1n,1-\frac1n\Bigr],
\]
and a real number in $[1/n,1-1/n]$ is at distance at least $1/n$ from every integer.
\end{proof}

The constant is attained, and for a reason rather than by coincidence.

\begin{proposition}[Computed; lower bound proved for all odd $n\ge5$; equality computed for $n\le41$]\label{prop:attained}
For every odd $n\ge5$, $\ \ell\bigl(n,\{(n-1)/2\}\bigr)\ \ge\ \dfrac{1}{n(n-1)}$, with equality for
every odd $n$ with $5\le n\le41$.
\end{proposition}

\begin{proof}
Put $b=(n-1)/2$, so $s=n-b=(n+1)/2$. Since $s$ and $b$ differ by $1$ they are coprime, so the least
$j\ge s$ with $\gcd(j,b)=1$ is $j_0=s$ itself. Also $2b=n-1<n$, so the speed $2b$ is present and
Proposition~\ref{prop:missing} applies, giving
\[
  \ell\bigl(n,\{b\}\bigr)\ \ge\ \frac{2s-j_0}{2nb\,j_0}\ =\ \frac{s}{2nb\,s}\ =\ \frac{1}{2nb}
  \ =\ \frac{1}{n(n-1)} .
\]
Note also that $b$ is coprime to $n$, since any common divisor of $(n-1)/2$ and $n$ divides
$\gcd(n-1,n)=1$, so Theorem~\ref{thm:unit} gives the same lower bound independently.

The matching upper bound is not proved here. It is verified by exact rational computation of
$\ell(n,\{b\})$ for every odd $n$ with $5\le n\le41$, where equality holds in every case; the
checker is committed. Proposition~\ref{prop:ceiling} gives only $\ell\le(n-2)/(n(n-1))$, which is
larger by a factor $n-2$, so the upper bound needs an argument the present method does not
supply.
\end{proof}

Neither primality of $n$ nor $n$ being large plays any role. When $n$ \emph{is} prime, more is true:
$b=(n-1)/2$ is then the worst single deletion there is. Over all $76$ primes with $5\le n\le400$ the
minimum of $\ell(n,\{b\})$ over $b$ equals $1/(n(n-1))$ and is attained at $b=(n-1)/2$, in every
case; and $1/(n(n-1))$ bounds $\ell(n,\{b\})$ from below for every $b$, which is
Theorem~\ref{thm:unit} again since every $b<n$ is coprime to a prime $n$. So at prime $n$ the
covering criterion of Theorem~\ref{thm:uncond} can never fire at the extremal deletion: it would
need $1/(n(n-1))>2/(n(n-2))$, that is $0>n$.

For even $n$ the extremal deletion is not so simply described. The value at $b=n-2$ has the closed
form $\ell(n,\{n-2\})=\frac{2}{n(n-2)}\bigl(\frac1P+\frac1{n-2-P}\bigr)$, with $P$ the least prime
not dividing $n-2$, exact for all even $6\le n\le400$ except $n=6,8$; at $n=32$ this reads
$\frac{2}{960}(\frac17+\frac1{23})=\frac1{2576}$. But $b=n-2$ is not always the minimiser: over the
$148$ even $n$ with $6\le n\le300$ it is so in $121$ cases and not in $27$, the first exceptions
being $n=10,14,18,24$.

\section{Computations}\label{sec:comp}

All arithmetic is exact over $\Z$ and $\mathbb{Q}$ in \texttt{i128}; no floating point is
load-bearing. The programs are \texttt{code/perturbation/perturb\_*.rs} in the accompanying archive,
the subcommands and the values each regenerates are documented in
\texttt{code/perturbation/README.md}. Components of $U(n,M)$ are obtained by sorting the endpoints
$(in\pm1)/(jn)$ of the excluded intervals and sweeping. Each range is stated where it is used; they
are not uniform.

\subsection*{A complete decision at \texorpdfstring{$k=2$}{k=2}}

For $5\le n\le100$ and every $M$ with $|M|=2$ there is no admissible $R$ at all: the conjecture holds
for every $2$-speed modification in that range. By Remark~\ref{rem:general} that is exactly the
statement that every set of $n-1$ distinct positive integers with precisely two members at least $n$
is lonely, for those $n$.

This is a decision rather than a sample, because two rigorous bounds make the enumeration finite.
Theorem~\ref{thm:cover} gives $\min R\le 4/(\ell(n,M)(n-4))$, which bounds the smaller inserted
speed; over $n\le30$ the largest value taken is $136$. Given $v_1$, the residual
$U(n,M)\setminus B_{v_1}$ must lie inside $B_{v_2}$, and $B_{v_2}$ is a \emph{disjoint} union of open
intervals, so each connected piece of the residual lies in a single one of them: a residual
component of length $L$ forces $v_2<2/(Ln)$. If the residual has no component of positive length it
is a finite set of rationals $p/q$, and $\nrm{v_2p/q}<1/n$ depends only on $v_2$ modulo $q$, so one
period of $\operatorname{lcm}(q)$ decides that branch; it did not arise.

Totals over $5\le n\le100$: $161\,696$ deletion sets, $7\,520\,944$ pairs $(M,v_1)$ examined and
$13\,989\,171$ enumeration steps,
no covering $R$, and no branch truncated by a cap. The conjecture is known for at most $13$
runners: seven in \cite{bs}, eight in \cite{rosenfeld}, nine and ten in \cite{trak}, and eleven to
thirteen in \cite{st}. So $161\,414$ of those $161\,696$ sets lie beyond the range proved outright.

The enumeration does not carry over to $k=3$ unchanged, and we record why. The second of the two
bounds above is a connectivity argument: it reads off $v_2<2/(Ln)$ from the fact that a residual
component must lie inside a single interval of $B_{v_2}$. That inference needs exactly one speed to
remain. With two still to place, a residual component may be covered jointly by
$B_{v_2}\cup B_{v_3}$ and neither speed need satisfy the bound. At $k=3$ the recursion reaches a
level with two speeds remaining, so the connectivity bound is outside its hypothesis there. The
repair is to use Theorem~\ref{thm:cover} on the residual at every level where $r\ge2$ speeds remain,
giving $\min\le 2r/(\ell(n-2r))$, which is uniform in $r$ and is the same measure argument that
removes the cap at level zero.

With that bound in place the enumeration at $k=3$ decides $7\le n\le40$: $3\,021\,814$ pairs
$(M,v_1)$ and $153\,363\,483$ enumeration steps, no covering $R$, and no branch truncated by a cap.
The corrected bound multiplies the number of steps by about seven relative to the connectivity
version, which is the measure of how much of the space the latter was missing. The $k=2$ decision is
unaffected either way, since exactly one speed remains at every recursive level there.

\paragraph{Reproduction status.} The $k=2$ totals above are regenerated by \texttt{code/perturbation/perturb\_coverk.rs 2 5 100}: the rows $5\le n\le100$ sum to $161\,696$ sets, $7\,520\,944$ pairs and $13\,989\,171$ steps with no cover and no capped branch (rows $31$ to $66$ come from a run of 17 October 2026, the others from 18 October 2026, same program). The $k=3$ totals are the column sums of the log of the original run (\texttt{data/repro\_logs/perturb\_coverk\_k3\_7\_40\_original.log}); the archived program reproduced its rows for $7\le n\le22$ exactly, and the rows $23\le n\le40$ were not re-run. The $k=3$ decision for $23\le n\le40$ is therefore a COMPUTED result of a development run that is not regenerated here.

Zhang \cite{zhang} reports exact $k=3$ searches inside windows (removals in the mid-range, $n\le140$; see his Remark~6.3 and Section~8); we have not found a height-free decision for $k=3$ in the literature read, and have not searched exhaustively. The case $k=1$ is settled by \cite[Thm.~2.3]{gw} together with \cite[Thm.~1.2, Lemma~2.3]{zhang}
and a computation; \cite[Thm.~2.4]{gw} gives finiteness for fixed $r$ but is not effective, and
\cite[Thm.~1.4]{zhang} computes the components in a sharper form than anything here. For $k=2$,
\cite[Thms.~7.5, 7.10, 7.13]{zhang} prove non-existence of tight two-speed modifications for the removed
sets $\{2,r'\}$ ($n\ge19$) and $\{q,r'\}$ ($n\ge\max(40,10q)$), and give a census for $8\le n\le140$
with inserted speeds at most $8n$; version~1 had verified $6\le n\le34$ under $\max R\le4n$. The
computation here is independent of those arguments: Theorem~\ref{thm:cover} bounds $\min R$ with no
hypothesis on the inserted speeds, so the enumeration is finite without a cap, and it covers
$5\le n\le100$ for $k=2$ and $7\le n\le40$ for $k=3$. For $k=2$ it is not a stronger result than
Zhang's theorems on the removed sets they cover.

As a control, the tight set $\{1,4,5,6,7,11,13\}$ at $n=8$, due to Goddyn and Wong \cite{gw} and
also recorded by Zhang, which is the $2$-speed modification with $M=\{2,3\}$ and $R=\{11,13\}$, lies
in the range swept here. It is tight, not a counterexample,
and the search correctly reports no covering: the times $1/8$ and $7/8$ survive, since neither $11$
nor $13$ is divisible by $8$.

The search is instrumented against the failure mode of a test that cannot succeed. Six positive
controls check that the interval subtraction removes a contained interval, trims a partial overlap to
the exact expected endpoint, leaves a disjoint interval untouched, splits a straddling interval in
two, and that the congruence branch fires at $t=1/4$ and $t=1/2$, where a covering residue class
exists. The negative control is that the search reports nothing for $n\le13$, where the conjecture is
a theorem.


\subsection{The tight $2$-speed modifications, exhaustively}\label{sec:tight2}

The same enumeration, run to decide tightness rather than to look for counterexamples, classifies
the tight $2$-speed modifications outright. With no cap on the inserted speeds, over
$1.41\cdot10^{8}$ enumeration steps and with no branch left undecided:

\begin{proposition}[Computed]\label{prop:tight2}
For $5\le n\le100$, the following three $2$-speed modifications of $[n-1]$ are tight, and the
enumeration finds no others:
\[
  \{2,4,6,8\}\ (n=5),\qquad \{1,4,5,6,7,11,13\}\ (n=8),\qquad
  \bigl([73]\setminus\{70,72\}\bigr)\cup\{140,144\}\ (n=74).
\]
\end{proposition}

Two caveats on completeness, both in the same direction. The search records, for each deleted pair
and each smallest inserted speed, the first covering completion it finds rather than all of them, so
a second tight set sharing both could in principle go unlisted; enumerating all completions instead
changes nothing over $5\le n\le34$, where the raw hits collapse to the same two sets. And the
interval arithmetic represents the complement of a closed bad set by closed intervals, so it retains
endpoints touched exactly and thereby over-estimates the surviving region. Both errors can only
cause a covering to be missed, never invented, so Proposition~\ref{prop:tight2} should be read as
listing the tight modifications the enumeration finds, with no counterexample to the conjecture
anywhere in the range, rather than as a proof that no fourth tight set exists.

None of the three is new, and it is worth being precise about why. The first is the dilation
$2\cdot[4]$. The second is the Goddyn--Wong set above. The third is an instance of the
simultaneous-acceleration theorem of Goddyn and Wong \cite{gw}, stated as \cite[Thm.~12]{ps}:
accelerating $r$ by a factor $m$ within $[n-1]$ preserves tightness when $r$
shares a factor with every element of $[n-r,\,m(n-r)-1]$, and accelerations at several $r$ compose.
Here $r=70$ and $r=72$ are both accelerated by $m=2$, with intervals $[4,7]$ and $[2,3]$; since
$70=2\cdot5\cdot7$ meets each of $4,5,6,7$ and $72$ meets each of $2,3$, both qualify.

So the content of Proposition~\ref{prop:tight2} is negative, and that is its interest. A complete
characterisation of tight instances is open, and the converse of that theorem is known to fail in
general \cite{ps}. Within $2$-speed modifications, up to $n=100$, the converse holds: every tight
instance is a dilation or a Goddyn--Wong acceleration, with no sporadic example.

Tightness of the three was re-checked independently of the search, by computing $\LR(V)$ in exact
integer arithmetic over the candidate times $a/q$ with $q$ dividing some $|v_i\pm v_j|$ or $2v_i$.
That routine returns $\LR=1/4,1/5,1/6$ on $\{1,2,3\}$, $\{1,2,3,4\}$, $\{1,3,4,5,9\}$, and returns
$2/13$ on a perturbation of the $n=8$ set, so it distinguishes tight from non-tight rather than
reporting tightness unconditionally.

\begin{remark}[The closest unconditional result on this object]\label{rem:bp}
For large $n$, Bohman and Peng \cite[Thm.~1.3]{bp}, with the range improved by Pomerance
\cite{pomerance}, show that a set with $n<\max R\le 2n-c(\log n)^{2}$ is not tight. In the notation
here that settles every $k$-speed modification, for every $k$, whose inserted speeds all lie below
$2n-c(\log n)^{2}$; their witness time is again a fraction whose denominator is assembled from
deleted speeds. It is asymptotic and ineffective in $n$, so it does not decide any particular small
$n$, which is what Section~\ref{sec:comp} does; but it is the strongest existing unconditional
statement about the objects of this paper.
\end{remark}

\section{The divisibility obstruction}\label{sec:divis}

Covering a rational point of the uncovered region is not an analytic condition but a divisibility
one, and the two families of rational points that always lie in $U(n,M)$ impose competing demands.

This section is expository: the content is known. Proposition~\ref{prop:divbase} is the mechanism
behind \cite[Lemma~2.3]{zhang}, which is in turn immediate from \cite[Lemma~2.2]{gw}, and
Proposition~\ref{prop:divtriv} is the same observation at the base point $1/n$. We include the
statements because the Lean development uses them and because they are what the search of
Section~\ref{sec:comp} prunes on, not as new results.

\begin{proposition}[Proved; Divisibility at a base point]\label{prop:divbase}
Let $b<n$, $\gcd(a,b)=1$, and let $v$ satisfy $\nrm{v\,a/b}<1/n$. Then $b\mid v$.
\end{proposition}

\begin{proof}
Write $r=va\bmod b$, so $\nrm{v\,a/b}=\min(r,b-r)/b$. If $r\ne0$ this is at least $1/b>1/n$. Hence
$r=0$, that is $b\mid va$, and $\gcd(a,b)=1$ gives $b\mid v$.
\end{proof}

\begin{proposition}[Proved; The times $j/n$]\label{prop:divtriv}
Let $\gcd(j,n)=1$. Then $j/n\in U(n,M)$ for \emph{every} $M\subseteq[n-1]$, and any $v$ with
$\nrm{v\,j/n}<1/n$ satisfies $n\mid v$.
\end{proposition}

\begin{proof}
A present speed $q$ has $1\le q<n$, so $n\nmid qj$ and $\nrm{q\,j/n}\ge1/n$; this uses nothing about
$M$. For the second claim, $\nrm{v\,j/n}<1/n$ forces $vj\equiv0\pmod n$, and $\gcd(j,n)=1$ gives
$n\mid v$.
\end{proof}

Both are equivalences: $b\mid v$ makes $a/b$ a centre of a bad interval of $v$, and likewise for
$n\mid v$ and $j/n$. We verified them as such in exact arithmetic, by
\texttt{code/perturbation/perturb\_divis.rs}: $593\,442$ instances of Proposition~\ref{prop:divbase} over
$3\le n\le40$, $1\le b<n$, $a$ a unit mod $b$ and $1\le v\le 3n$; and $132\,684$ instances of
Proposition~\ref{prop:divtriv} over $3\le n\le60$ with $1\le v\le 3n$. No exception. The membership
$j/n\in U(n,M)$ was checked over \emph{all} $2^{n-1}$ subsets $M$ for $3\le n\le14$, including
$M=[n-1]$: $121\,368$ instances, no exception.

\begin{corollary}[Proved; Necessary conditions on $R$; cf.\ \protect{\cite[Lemmas~2.3, 6.1]{zhang}}]\label{cor:divnec}
Suppose $\LR(([n-1]\setminus M)\cup R)<1/n$. Then some $v\in R$ is divisible by $n$; and for every
$b\in M$ having no present multiple, some $v\in R$ is divisible by $b$. Writing $M_0$ for the set of
such $b$, if $|M_0|=k$ then the $k$ speeds of $R$ must realise the $k+1$ demands $M_0\cup\{n\}$, so
some $v\in R$ is divisible by $\operatorname{lcm}(d,d')$ for two distinct demands $d,d'$.
\end{corollary}

The statement differs from \cite[Lemma~2.3]{zhang} in three ways: he assumes $2r>n-1$ where the
condition here is the weaker ``no present multiple'' (which $2b\in M$ also satisfies), he assumes
tightness where we assume $\LR<1/n$, and the clause about $n\mid v$ and the pigeonhole over
$M_0\cup\{n\}$ are not in the statements we read (his Lemma~6.1 extends the divisibility to every
removed $r<n$ under tightness). The mechanism is his.

\begin{proof}
If $b$ has no present multiple then $\delta^{0}=0$, so $a/b\in U(n,M)$ by Theorem~\ref{thm:exact};
$j/n\in U(n,M)$ always by Proposition~\ref{prop:divtriv}. Each must be covered by some $v\in R$, and
Propositions~\ref{prop:divbase} and~\ref{prop:divtriv} identify the divisor. The last clause is
pigeonhole on $k+1$ demands and $k$ speeds.
\end{proof}

At $k=1$ the two demands collide on a single speed, and that alone settles the regime.

\begin{theorem}[Proved; Criterion D]\label{thm:critD}
Let $n\ge3$, $k=1$, and $M=\{b\}$ with $2b\ge n$, so that $b$ has no present multiple. If
\[
  \ell(n,\{b\})\cdot n\cdot\operatorname{lcm}(b,n)\ \ge\ 2,
\]
then $\LR(([n-1]\setminus\{b\})\cup R)\ge1/n$ for every admissible $R$.
\end{theorem}

\begin{proof}
Suppose not, and let $R=\{v\}$. By Corollary~\ref{cor:divnec}, $b\mid v$ and $n\mid v$, so
$\operatorname{lcm}(b,n)\mid v$ and in particular $v\ge\operatorname{lcm}(b,n)$. The longest
component $W$ of $U(n,M)$ is connected and contained in $B_v$, which is a disjoint union of open
intervals, so $W$ lies in one of them and $\ell(n,\{b\})=|W|<2/(vn)$. Hence
$v<2/(\ell n)$, contradicting the displayed hypothesis.
\end{proof}

Theorem~\ref{thm:critD} is not new, and we claim no priority for it. The regime $2b\ge n$ at
$k=1$ is already closed in the literature by exactly this route: \cite[Lemma~2.3]{zhang} gives
$b\mid v$, after which \cite[Thms.~1.1, 1.2]{zhang} classify the single-speed modifications. We record the criterion in this form only because it
is what the formalization proves and because the numerical comparison below locates the measure
criterion's loss. Exhaustively over all $b$ with $2b\ge n$, Theorem~\ref{thm:critD} certifies
$100\%$ of them at $n=20,26,32,38,44,50$, against $30.0,30.8,37.5,31.6,40.9,40.0$ per cent for the
criterion of Theorem~\ref{thm:uncond}, and $100\%$ is also the truth. So on this regime the measure
criterion is losing roughly $60$ per cent of what divisibility sees, which is the useful reading.
The mechanism is visible at $n=32$, $M=\{30\}$: the covering bound caps $v$ at $171$, while
$\operatorname{lcm}(30,32)=480$.

For $k\ge2$ the two demands need not fall on the same speed, and Corollary~\ref{cor:divnec} becomes
a constraint rather than a contradiction: the speed covering the base point of the longest window is
divisible by $b$, hence at least $b\lceil n/b\rceil$, and so covers only the radius $1/(vn)$ about
$a/b$, leaving two arms to the remaining $k-1$ speeds. Each arm is connected, so the connectivity
bound applies to it, and the analysis is the finite recursion of Section~\ref{sec:comp} rather than
a closed criterion. What the divisibility does supply there is a sound and substantial pruning: by
Proposition~\ref{prop:divtriv}, if any $j/n$ survives the speeds chosen so far, the last remaining
speed must be a multiple of $n$.

Propositions~\ref{prop:divbase} and~\ref{prop:divtriv}, the membership $j/n\in U(n,M)$, and the
divisibility half of Theorem~\ref{thm:critD} are formalized; see Section~\ref{sec:comp} and the
table in the Formalization section.

\begin{itemize}
\item Theorem~\ref{thm:unit} was tested on all $M$ of size $1$ with $n\le 44$ ($602$ instances with
$M$ meeting the units mod $n$) and of size $2$ with $n\le34$ ($4857$ instances): no violation, and
the value $1/(n(n-1))$ attained.
\item The lower bound of Proposition~\ref{prop:missing}, together with
Theorem~\ref{thm:unit}, was compared against the exhaustively computed $\ell(n,\{b\})$ for all
$4004$ pairs $b<n\le 90$: no violation. Over that range
$\min \ell(n,\{b\})\,n(n-1)P(b)=1.59563$, at $(n,b)=(74,70)$, and the same minimum for the proved
lower bound is $1.39048$, at the same pair.
\item Over all $b<n\le400$: the largest coprime gap $j_0-s$ is $9$, at $(n,b)=(212,210)$; and in
the regime where it is used, namely $2b\le n-1$, the quantity $(2s-j_0)/s$ is always at least
$0.625$, attained at $(n,b)=(14,6)$, with no case where $2s-j_0\le0$.
\item $P(b)<2\log b+2$ was verified for all $2\le b\le 2\cdot10^{5}$, consistent with the case check
in the proof of Theorem~\ref{thm:uniform}(ii).
\end{itemize}

The logarithm in Theorem~\ref{thm:uniform}(iii) is not an artifact of the proof, but establishing
that needs an \emph{upper} bound on $\ell$, which Proposition~\ref{prop:missing} does not provide.
Zhang \cite[Thm.~1.4]{zhang} computes $\ell(n,\{b\})$ exactly; his formula at $M=\{n-2\}$ gives
$\ell\,n(n-1)=2/P(n-2)+o(1)$ as $n\to\infty$ along any sequence with $P(n-2)$ fixed, and
$P(b)\to\infty$ along primorials, so no bound of the form $c/n^{2}$ holds for all $M$. We record the
agreement of our computed values with that formula, without relying on it elsewhere: $P=7$ for $b=720$
gives limit $2/7=0.2857$ and measured $0.28892$ at $n=722$; $P=11$ for $b=420$ gives $0.18182$ and
measured $0.18715$ at $n=422$; $P=13$ for $b=2310$ gives $0.15385$ and measured $0.15478$ at
$n=2312$.

\section{Formalization}

The covering bound, the coprime component bound, the localization theorem and the unconditional
criterion are formalized in Lean~4 with Mathlib. Each declaration builds with no \texttt{sorry} and
depends only on \texttt{propext}, \texttt{Classical.choice} and \texttt{Quot.sound}. The development
comprises $70$ theorems and $5$ definitions, across the fifteen modules
\texttt{Perturb*.lean}, \texttt{CoveringBound.lean} and \texttt{PhaseBound.lean}; the count is
reproduced by
\texttt{grep -cE '\textasciicircum(theorem|def)' } on those files.

\begin{center}
\begin{tabular}{lll}
\hline
Statement & Lean declaration & constant \\
\hline
Cor.~\ref{cor:main}, coprime case & \texttt{perturbation\_bound\_of\_coprime} & $(2{+}4/n)k$ \\
Thm.~\ref{thm:cover}, assembled & \texttt{covering\_bound} & $(2{+}4/n)k$ \\
Lemma~\ref{lem:strip} & \texttt{volume\_badSet\_inter\_Icc} & $(2{+}4/n)/(vn)$ \\
$B_v$ as a union of intervals & \texttt{badSet\_eq} & n/a \\
Thm.~\ref{thm:cover}, arithmetic core & \texttt{covering\_forces\_small\_speed} & general \\
Thm.~\ref{thm:unit} & \texttt{perturb\_component\_of\_coprime} & sharp \\
Thm.~\ref{thm:unit}, core & \texttt{perturb\_component} & sharp \\
uniqueness of $j_+$ & \texttt{residue\_one\_unique} & n/a \\
$r\ne0$ & \texttt{residue\_ne\_zero} & n/a \\
distance to $\Z$ on $[1/n,1-1/n]$ & \texttt{dist\_int\_ge\_of\_mem\_Icc} & n/a \\
Lemma~\ref{lem:mono} & \texttt{uncovered\_mono} & n/a \\
Thm.~\ref{thm:unit} in region form & \texttt{Icc\_subset\_uncovered} & n/a \\
Prop.~\ref{prop:missing}, containment & \texttt{Icc\_subset\_uncovered\_missing} & sharp \\
Prop.~\ref{prop:missing}, upper edge & \texttt{hnon\_of\_j0} & sharp \\
Cor.~\ref{cor:main}, missing-speed case & \texttt{perturbation\_bound\_missing} & $(4{+}8/n)k$ \\
maximality of $j_0$ & \texttt{le\_of\_congr\_of\_ge} & n/a \\
$j_0$ exists in the window & \texttt{exists\_coprime\_window} & n/a \\
existence of $a$ & \texttt{exists\_a\_of\_coprime} & n/a \\
maximality, packaged & \texttt{hj0max\_of\_spec} & n/a \\
Thm.~\ref{thm:uniform}(i), two regimes & \texttt{exists\_interval\_missing} & explicit \\
\textbf{Thm.~\ref{thm:uniform}(i)} & \texttt{exists\_interval\_uncovered} & explicit \\
Thm.~\ref{thm:uncond}, covering step & \texttt{not\_covered\_of\_long\_interval} & $(2{+}4/n)k$ \\
\textbf{Thm.~\ref{thm:uncond}} & \texttt{exists\_lonely\_time} & $(2{+}4/n)k$ \\
Cor.~\ref{cor:uncond}, window from $j_0$ & \texttt{Icc\_subset\_uncovered\_of\_j0} & sharp \\
Cor.~\ref{cor:uncond}, criterion & \texttt{lonely\_of\_j0\_criterion} & $(2{+}4/n)k$ \\
Cor.~\ref{cor:uncond}(a) & \texttt{lonely\_of\_one\_mem} & $n>(6{+}8/n)k$ \\
\textbf{Cor.~\ref{cor:uncond}(b)} & \texttt{lonely\_of\_criterion\_b} & $(4{+}8/n)kbj_0$ \\
Cor.~\ref{cor:uncond}(c), one-sided & \texttt{lonely\_of\_criterion\_c} & vacuous \\
Prop.~\ref{prop:twosided}, left edge & \texttt{missing\_nonmultiple\_protects\_left} & sharp \\
\textbf{Prop.~\ref{prop:twosided}} & \texttt{Icc\_subset\_uncovered\_twosided} & sharp \\
\textbf{Cor.~\ref{cor:uncond}(c)} & \texttt{lonely\_of\_criterion\_c'} & $(2{+}4/n)kbj_0j_1$ \\
\textbf{Thm.~\ref{thm:localize}} & \texttt{exists\_deleted\_fraction} & sharp \\
Thm.~\ref{thm:localize}, denominator form & \texttt{exists\_deleted\_denominator} & sharp \\
Thm.~\ref{thm:localize}, lowest terms & \texttt{deleted\_denominator\_of\_dvd} & sharp \\
Cor.~\ref{cor:measure}, covering step & \texttt{uncoveredStrict\_subset\_union} & sharp \\
\textbf{Thm.~\ref{thm:exact}}, abutment & \texttt{uncovered\_between} & sharp \\
uncovered $\to$ strict on an interval & \texttt{uncoveredStrict\_of\_Icc\_subset} & sharp \\
\textbf{Thm.~\ref{thm:exact}}, abutment left & \texttt{uncovered\_between\_left} & sharp \\
two-sided window, any regime & \texttt{Icc\_subset\_uncovered\_twosided\_general} & sharp \\
two-sided containment, per speed & \texttt{Icc\_subset\_uncovered\_twosided\_direct} & sharp \\
\textbf{Cor.~\ref{cor:window}}, two-sided & \texttt{lonely\_of\_window\_twosided} & $(2{+}4/n)k$ \\
Cor.~\ref{cor:measure} on a component & \texttt{exists\_base\_point\_of\_Icc} & sharp \\
\textbf{Cor.~\ref{cor:window}} & \texttt{lonely\_of\_window} & $(2{+}4/n)k$ \\
\textbf{Prop.~\ref{prop:divbase}} & \texttt{dvd\_of\_covers\_base} & sharp \\
Prop.~\ref{prop:divtriv}, membership & \texttt{mem\_uncovered\_of\_coprime} & sharp \\
\textbf{Prop.~\ref{prop:divtriv}} & \texttt{dvd\_of\_covers\_trivial} & sharp \\
\textbf{Thm.~\ref{thm:critD}}, divisibility half & \texttt{exists\_survivor\_of\_not\_lcm\_dvd} & sharp \\
\hline
\end{tabular}
\end{center}

What is and is not machine-checked. The list below is the single authoritative inventory; the
abstract and the table agree with it, and the table lists every declaration referred to here.

\begin{itemize}
\item \textbf{Formalized:} Lemma~\ref{lem:strip}; Theorem~\ref{thm:cover}; Lemma~\ref{lem:mono};
Theorem~\ref{thm:unit}; Proposition~\ref{prop:missing}; Proposition~\ref{prop:twosided}; the
selection of $j_0$ and of $a$; Theorem~\ref{thm:uniform}(i); Theorem~\ref{thm:uncond};
Theorem~\ref{thm:localize}; Propositions~\ref{prop:divbase} and~\ref{prop:divtriv} in full, and
the divisibility half of Theorem~\ref{thm:critD}; the abutment half of Theorem~\ref{thm:exact} and its containment
half in the restricted form described below; Corollary~\ref{cor:window};
Corollary~\ref{cor:uncond} in all three cases; and \emph{both} displays of
Corollary~\ref{cor:main}.
\item \textbf{With a weaker constant.} The formalized strip lemma is proved by a union bound over
the translates meeting the window rather than by periodicity. The translates that can meet
$[x,x+\ell]$ are indexed by the integers in an open interval of length $v\ell+2/n$, of which there
are at most $v\ell+1+2/n$, and each contributes at most $2/(vn)$; so Lean proves
\[
  \lambda\bigl(B_v\cap[x,x+\ell]\bigr)\ \le\ \frac{2\ell}{n}+\frac{2+4/n}{vn}
\]
in place of the sharp $2\ell/n+2/(vn)$. Everything downstream inherits the constant $2+4/n$ rather
than $2$. What Lean proves is the second display of Theorem~\ref{thm:cover}, the bound on
$\min R$; the first display, the inequality for $\sum_{v\in R}1/v$, appears in the development only
as a hypothesis of \texttt{covering\_forces\_small\_speed} and is not itself formalized. The two
displays of Corollary~\ref{cor:main} are formalized in the forms $\min R\le(4{+}8/n)k\,n^{3}/(n-2k)$
and
$\min R\le(2{+}4/n)k\,n(n-1)/(n-2k)$. Theorem~\ref{thm:uncond} is formalized with $(2+4/n)k$ in
place of $2k$; the loss is a factor $(n+2)/n$, and reaching $2$ exactly would need the periodicity
argument.
\item \textbf{What that costs in reach.} The machine-checked form of Corollary~\ref{cor:window}
requires $w(n,b,a,M)>(2n+4)k/(n^{2}(n-2k))$ rather than $2k/(n(n-2k))$, and it is not the window
$w$ of Definition~\ref{def:window} that \texttt{lonely\_of\_window} consumes but a smaller one: its
hypothesis $\delta_0\ge0$ makes the window one-sided, so it cannot straddle the base point. Its
offset and far cap are now the sharp ones, $\delta_0\ge1/(jn)$ and $\delta_1\le1/n$; the crude
$1/(2bn)$ and $1/(nb)$ survive only one level up, in \texttt{Icc\_subset\_uncovered\_of\_j0}, and it
is there that the $4358$ base points with $\delta^{+}>1/(nb)$, and the $488$ at which $b$ has a
present multiple and $\delta^{0}<1/(2bn)$, apply. The figures below are for the union of \texttt{lonely\_of\_window},
\texttt{lonely\_of\_criterion\_c'} and \texttt{lonely\_of\_one\_mem}. Counted exhaustively
over all $k$-subsets at $n=22,26,30,34,38,42$, what Lean certifies is
$38.10,32.00,37.93,42.42,35.14,39.02$ per cent at $k=1$; $11.90,18.33,16.50,20.83,19.52,21.83$ at
$k=2$; and $14.29,12.17,11.82,17.80,16.74,15.43$ at $k=3$, against the criterion's own
$38.10,36.00,41.38,42.42,37.84,41.46$; $21.43,19.33,23.89,21.78,24.77,22.44$; and
$14.44,13.13,11.88,18.37,16.82,15.55$.
\item \textbf{Not assembled in Lean, and formalized only in part.} Three statements of this paper
are not fully machine-checked, and we set out exactly what is missing.
Theorem~\ref{thm:exact}: both abutment halves are formalized, \texttt{uncovered\_between} to the
right and \texttt{uncovered\_between\_left} to the left, and
\texttt{Icc\_subset\_uncovered\_twosided\_general} assembles them into a two-sided containment with
no restriction on the regime, the hypothesis being that $b$ has no present multiple rather than
$n\le2b$. The containment \texttt{Icc\_subset\_uncovered\_missing} now carries the sharp per-speed offset
$\delta_0\ge1/(jn)$ in place of $1/(2bn)$, and the far cap $\delta_1\le1/n$ in place of $1/(nb)$;
\texttt{Icc\_subset\_uncovered\_twosided\_direct} gives the two-sided interval from per-speed
conditions alone whenever $b$ has no present multiple, and \texttt{lonely\_of\_window\_twosided} is
the resulting criterion. Together these lift the machine-checked reach at $k=2$ from
$10.00,15.67,14.04,17.61,16.07,18.78$ to the figures above. What remains unformalized on
Theorem~\ref{thm:exact} is the maximum over base points.
Corollary~\ref{cor:measure}: the covering statement is formalized, the summation
$\sum_{q\in M}q\cdot 2/(qn)=2|M|/n$ is not. Theorem~\ref{thm:cover}: the first display is not
formalized, as noted above. Theorem~\ref{thm:critD}: the divisibility half is
\texttt{exists\_survivor\_of\_not\_lcm\_dvd}, in full; the numerical assembly against the covering
bound is not formed in Lean.
\item \textbf{Not formalized:} parts (ii) and (iii) of Theorem~\ref{thm:uniform}, which rest on
Chebyshev-type input and on Iwaniec's theorem respectively; Proposition~\ref{prop:ceiling}; the
elementary implication of Remark~\ref{rem:closed}; the emptiness computation of
Remark~\ref{rem:cempty}; and the recursion of Section~\ref{sec:open}. Remark~\ref{rem:zhang}
records a comparison and carries no formal content.
\end{itemize}

The row \texttt{lonely\_of\_criterion\_c} records the superseded one-sided form of case (c) and is
provably vacuous; \texttt{lonely\_of\_criterion\_c'} replaces it and is not, though by
Remark~\ref{rem:cempty} it fires only at $k=1$.

The Lean development depends on no external result beyond Mathlib. In particular
Proposition~\ref{prop:missing} is proved there from scratch, Theorem~\ref{thm:localize} uses only
Mathlib's \texttt{Real.exists\_nat\_abs\_mul\_sub\_round\_le}, and Remark~\ref{rem:zhang} is used
nowhere.

\section{The covering bound away from modifications}\label{sec:beyond}

The covering bound was set up for modifications of $[n-1]$, but nothing in it refers to that
structure. Stated for an arbitrary family it reads as follows, and this form is machine-checked.

\begin{proposition}[Verified; Lean: \texttt{covered\_interval\_short}]\label{prop:short}
Let $v>0$ and $0<\delta\le\tfrac12$. If $\nrm{vt}<\delta$ for every $t\in[x,y]$, then
$(y-x)v<2\delta$.
\end{proposition}

The times at which a single speed fails to separate form intervals of length exactly $2\delta/v$,
and an interval inside that union lies in one of them. The proof is the intermediate value theorem:
otherwise $vt$ takes the value $k+\delta$ somewhere in $[x,y]$, where the distance to the integers is
exactly $\delta$.

\begin{corollary}[Verified; Lean: \texttt{le\_ML\_of\_interval}]\label{cor:mlbound}
Let $V$ be a finite family of nonzero integers, $0<\delta\le\tfrac12$, and $j$ an index with
$V_j>0$. If some interval $[x,y]$ satisfies $\nrm{V_it}\ge\delta$ for every $i\ne j$ and every
$t\in[x,y]$, and $2\delta\le(y-x)V_j$, then $\LR(V)\ge\delta$.
\end{corollary}

We record one application, to the four-speed loneliness spectrum. Kravitz \cite{kravitz} conjectured
that $\LR$ of $n$ speeds is either $s/(ns+1)$ or at least $1/n$, proved it for $n\le3$, and proved it
for $n=4$ and $n=6$ under the lacunary hypothesis $v_n>4v_{n-1}^{4}$ \cite[Thm.~2.2]{kravitz}; the
conjecture at $n=4$ is false in general \cite{fansun}. The first gap
alone, that four speeds have $\LR=\tfrac15$ or $\LR\ge\tfrac29$, is open.

Write $\ell(v_1,v_2,v_3)$ for the longest interval on which all three of $v_1,v_2,v_3$ stay at
distance $\ge\tfrac29$. Corollary~\ref{cor:mlbound} applied with $\delta=\tfrac29$ and $j=4$ gives:
if $\LR(v_1,\dots,v_4)<\tfrac29$ then $v_4<\tfrac{4}{9}/\ell$. Computation in exact rational arithmetic over all coprime
triples with speeds at most $110$ gives $\min\ell\cdot v_3=\tfrac19$, attained only at $(1,2,3)$,
with second smallest value $\tfrac4{27}$, whence
$v_4<4v_3$, so the first gap holds throughout $v_4\ge4v_3$. That range is strictly wider than the
lacunary hypothesis $v_4>4v_3^{4}$ of \cite[Thm.~2.2]{kravitz}, since $4v_3^{4}$ exceeds $4v_3$ for
$v_3\ge2$; the conclusion, however, is weaker, settling only the first gap and not the full spectrum
statement. The minimum is a computation, not a theorem, so the constant $4$ is HEURISTIC here.

What the bound does not do is close the remaining range, and the obstruction is sharp rather than
technical. To force the contradiction $v_4<v_3$ one needs $\ell v_3>\tfrac49$, whereas the median of
$\ell v_3$ is $\tfrac59$ and the small values are $\tfrac19,\tfrac4{27},\tfrac29$. The triples
realising them are, with few exceptions, exactly those with $v_3=v_1+v_2$. On such a triple
$\nrm{tv_3}=\nrm{tv_1+tv_2}$, so with $u=tv_1$ and $w=tv_2$ the three conditions collapse to a single
condition on one point of $\mathbb{T}^2$, namely membership of
$\{(u,w):\nrm{u},\nrm{w},\nrm{u+w}\ge\delta\}$. At $\delta=\tfrac29$ that set is two congruent right
triangles of leg $1-3\delta=\tfrac13$, of total area $(1-3\delta)^2=\tfrac19$, the same constant
again. Deciding this case is a question about an orbit meeting a fixed planar region, which is the
subject of the main lemma of \cite{bhk}; we have not carried that out.

\subsection{A barrier for position-uniform covering lemmas}\label{sec:barrier}

The covering lemmas available for this kind of question all share a feature: their conclusion is
that a group meets \emph{every} rectangle of given dimensions, so the criterion they supply cannot
depend on where the rectangle sits. That is what makes them convenient, and it is also what limits
them here.

Fix $\delta=\tfrac29$ and consider four speeds with $v_3=v_1+v_2$. Fixing $tv_2\equiv\delta$ makes
the second condition automatic and reduces the remaining three to asking whether a translated
cyclic subgroup of $\mathbb{T}^2$ meets a rectangle. After normalising so that the subgroup is
generated by $(1/n,c/n)$, the question becomes whether $\{kc \bmod n\}$, as $k$ runs over an
interval of $L$ consecutive integers, meets an interval of length $M$; here $n,c,L,M$ are exactly
the data visible to a position-uniform criterion.

\begin{proposition}[Computed]\label{prop:barrier}
No criterion depending only on $(n,c,L,M)$ decides this question. Explicitly,
$(v_1,v_2,v_3,v_4)=(1,6,7,11)$ and $(1,6,7,17)$ both give $(n,c,L,M)=(6,5,2,3)$, and the required
$t$ exists for the second and not for the first. The pairs $(1,6,7,10)$ with $(1,6,7,22)$, and
$(1,7,8,15)$ with $(1,7,8,22)$, are further witnesses.
\end{proposition}

\begin{proof}
Each assertion is a finite check on the four listed pairs: the two members of a pair agree on
$(n,c,L,M)$ by direct substitution, and the existence or non-existence of the required $t$ is
decided by scanning the finitely many candidate times, since an optimal $t$ may be taken with
denominator a sum or difference of two speeds. A criterion depending only on $(n,c,L,M)$ assigns
the same verdict to both members of a pair and is therefore wrong on one of them.
\end{proof}

The consequence is that a position-uniform covering lemma cannot decide every instance of this
construction, whatever its constants. Measured on $147\,099$ instances, the best such criterion we
found, namely that the largest gap of the point set is smaller than $M$, decides $94.17\%$ of them
and never errs; the residue is where instances agreeing on $(n,c,L,M)$ disagree in outcome. This
accounts for a ceiling that four separate arguments met at the same place, and it says what a proof
must do instead: use the arithmetic of $v_1,v_2,v_4$, or a construction other than fixing
$tv_2\equiv\delta$.

\subsection{A Fourier criterion, and how it evades the barrier}\label{sec:fourier}

Proposition~\ref{prop:barrier} bounds what a position-uniform covering lemma can do. The following
criterion is not one, and is not subject to it.

Let $T\subseteq\R/\Z$ be a finite union of $N$ intervals of total measure $\mu$, let $v$ be a
positive integer, and write $F(k)=\int_T e^{-2\pi ikt}\,dt$. Say $t$ is \emph{good} if
$\nrm{vt}\ge\delta$.

\begin{proposition}[Proved; single mode]\label{prop:fourier1}
If no point of $T$ is good, then $\operatorname{Re}F(v)\ \ge\ \cos(2\pi\delta)\,\mu$.
\end{proposition}

\begin{proof}
If no point is good then $\nrm{vt}<\delta$ throughout $T$, so for each $t\in T$ the phase $2\pi vt$
lies within $2\pi\delta$ of a multiple of $2\pi$, whence $\cos(2\pi vt)\ge\cos(2\pi\delta)$ by
monotonicity of the cosine on $[0,\pi]$. Integrate over $T$. (Lean:
\texttt{cos\_lt\_cos\_of\_nid\_lt}, the pointwise step; the integration is not formalized.)
\end{proof}

Its contrapositive is a criterion: if $\operatorname{Re}F(v)<\cos(2\pi\delta)\mu$ then some point of
$T$ is good. Taking $T$ to be the uncovered set of the three slower speeds and $v=v_4$, this decides
$95.28\%$ of the instances left open at four speeds, against $70.33\%$ for the best criterion
obtainable from component lengths alone.

Using all modes rather than one sharpens it. With $g$ the indicator of the good phases,
\[
  \bigl|T\cap\{\nrm{vt}\ge\delta\}\bigr|
   \;=\;\hat g(0)\mu+2\operatorname{Re}\sum_{m\ge1}\hat g(m)\overline{F(mv)},
\]
and truncating at $M$ modes costs at most $2N/(\pi^2vM)$, since $|F(k)|\le N/(\pi|k|)$ for a union
of $N$ intervals and $|\hat g(m)|\le1/(\pi m)$. Positivity of the truncated sum less that tail
forces a good point. At $M=40$ this decides every one of $49\,121$ instances except
$(v_1,\dots,v_4)=(3,4,7,11)$, where $\LR=\tfrac29$ exactly, so the intersection has measure zero and
no criterion demanding positive measure can decide it.

The reason this passes the barrier is that $F$ depends on where $T$ sits, not only on its
translation-invariant data, so Proposition~\ref{prop:barrier} does not apply. What remains, and is
not proved here, is a lower bound for the truncated sum valid for every instance. Two structural
facts bear on it. First, $1_T$ is a product of three indicators, $1_T(t)=h(v_1t)h(v_2t)h(v_3t)$, so
its Fourier coefficients are a convolution
$F(k)=\sum_{av_1+bv_2+cv_3=k}\hat h(a)\hat h(b)\hat h(c)$, large only when $k$ is close to a small
integer combination of the slower speeds. Second, measurement of the margin across the same
$49\,121$ instances shows it never falls below $0.372\mu$ unless $v_4$ satisfies an integer relation
with $v_1,v_2$ having coefficients at most $3$; every instance with a small margin satisfies such a
relation, and those instances are rank two, where a different and elementary argument applies. So
the difficulty is confined to a case the criterion is not needed for. (The margin statement of this paragraph was measured in a session and its generating program is not archived; HEURISTIC. The counts $147\,099$, $94.17\%$, $49\,121$, $49\,120$ and the single undecided tuple $(3,4,7,11)$ are regenerated by \texttt{code/perturbation/g1\_gap\_criterion.rs} and \texttt{r1\_fourier\_criterion.rs}, which use floating point; $\LR(3,4,7,11)=2/9$ is exact, \texttt{code/obstructions/mlcheck.rs}.)

\section{Limits of the method}\label{sec:limits}

The covering bound is a measure estimate, and measure estimates on this problem are capped. We
record where the cap sits, because it bounds what any sharpening of Section~\ref{sec:uncond} can
achieve.

First, at the baseline the bad sets already have full measure. By Dirichlet every $t$ has some
$q\le n-1$ with $\nrm{qt}\le1/n$, so the set on which every $j\in[n-1]$ is protected is the finite
set $\{a/n:\gcd(a,n)=1\}$, of measure $0$. Hence $\lambda\bigl(\bigcup_{j\in[n-1]}B_j\bigr)=1$
exactly, which we verified in exact arithmetic for $n=5,7,11,13,17,23,31,41$. Any argument of the
shape ``the bad sets cannot cover because their measure is below $1$'' therefore fails at the
extremal case. This is the classical factor-of-two obstruction, recorded in \cite[Prop.~1.1]{tao}
and sharpened there by the example of the primes in $(n/4,n/2)$, for which the union bound is tight
to $1+O(1/n)$.

Second, in the modification setting the inserted speeds have a measure surplus. The $k$ speeds of
$R$ carry total bad-set measure exactly $2k/n$, while $\lambda(U(n,M))\le2k/n$ with the extremal
ratio $\lambda(U)\big/(2k/n)$ lying between $0.5139$ and $0.6492$ over all $M$ with $|M|\le3$ and
$n\in\{16,22,28,34,40\}$, and never reaching $1$. So $R$ always has between $35$ and $49$ per cent
more measure than the region it must cover, and no refinement of the measure step can force
non-covering.

Third, the criterion has an explicit ceiling. Its strength is governed by the smallest
$\ell(n,\{b\})$, and that minimum behaves like $2/(n^{2}P(n-2))$, with $P(m)$ the least prime not
dividing $m$: at $b=n-2$ one has $s=2$ and the binding unit is the least integer at least $2$
coprime to $n-2$. Over $n=50,200,800,3200$ the ratio of $\ell(n,\{n-2\})\,n^{2}$ to $2/P(n-2)$ is
$1.163,\,1.036,\,1.009,\,1.002$. The worst $n$ are those with $n-2$ a primorial, where
$\ell\,n^{2}\log n$ stays near $1$; this is why Theorem~\ref{thm:uniform}(ii) has the shape it does,
and it is sharp.

The conclusion we draw is that the results of this paper are the ones the method reaches, and that
progress on the conjecture itself will not come from sharpening the measure step. The modern
advances for small numbers of runners \cite{bs,rosenfeld,trak,st} are congruence sieves, a different
mechanism.

\section{Bounding every inserted speed}\label{sec:open}

Theorem~\ref{thm:cover} bounds $\sum_{v\in R}1/v$, hence $\min R$, but not every element of $R$. For
$k=1$ these coincide. For $k\ge2$ one would like to recurse: fix the smallest inserted speed $v_1$,
pass to $U(n,M)\setminus B_{v_1}$, and apply Theorem~\ref{thm:cover} again with $k-1$ speeds. That
needs a lower bound on the longest component of the residual region.

We tested the recursion with $v_1$ chosen adversarially to fragment the region, over all $M$ with
$|M|=2$, all $v_1\in[n,6n]$ and $8\le n\le34$, in exact rational arithmetic. In every instance the
residual region retained a component of positive length, and the minimum of
$\bigl(\text{longest component of } U(n,M)\setminus B_{v_1}\bigr)\cdot n(n-1)$ was $0.375$, attained
at $n=8$, $M=\{2,6\}$, $v_1=24$, with no downward trend.

\begin{proposition}[Proved; Ceiling]\label{prop:ceiling}
Let $n\ge3$. For every $v\ge1$, every connected subset of $[0,1)\setminus B_v$ has length at most
$(n-2)/(nv)$.
\end{proposition}

\begin{proof}
$[0,1)\setminus B_v$ is the disjoint union of the closed intervals
$\bigl[\tfrac{i+1/n}{v},\ \tfrac{i+1-1/n}{v}\bigr]$, $0\le i<v$, each of length $(1-2/n)/v$. A
connected set meeting two of them would meet the open gap between.
\end{proof}

There are thus two hard upper bounds on the longest component $\ell'$ of
$U(n,M)\setminus B_{v_1}$: the ceiling $(n-2)/(nv_1)$ of Proposition~\ref{prop:ceiling}, and
$\ell(n,M)$ itself, since the residual is a subset of $U(n,M)$. Neither dominates: for $v_1$ large the ceiling is the smaller, while at $v_1\approx n$ the ceiling
is about $1/n$ and $\ell(n,M)$ can be as small as order $1/n^{2}$, for instance at
$M=\{n/2,n-1\}$ with $n$ even. It is not always that small: $\ell(n,\{1,2\})=2/(3n)$ exactly, verified for $5\le n\le40$.

Consequently a conjecture asserting a fixed fraction of $\ell(n,M)$ alone, with no dependence on
$v_1$, is false: Proposition~\ref{prop:ceiling} forces $\ell'\to0$ as $v_1\to\infty$ while
$c\,\ell(n,M)$ stays fixed. A lower
bound independent of $v_1$ fails because the ceiling tends to $0$: at $n=8$, $M=\{2,6\}$ the longest
component is $6.696\cdot10^{-3}$ at $v_1=24$ and $1.144\cdot10^{-5}$ at $v_1=65536$, the latter
equal to the ceiling exactly, while the residual keeps positive measure: $0.031548$ at $v_1=24$ and
$0.054915$ at $v_1=65536$, the latter approximately the limiting value $\tfrac34\lambda(U(8,\{2,6\}))=123/2240=0.0549107$. Taking $v_1=n$ and minimising over $M$ with $|M|=2$, the ratio
$\ell'\big/\frac{n-2}{n^{2}}$ is $0.1000,0.0900,0.0653,0.0455,0.0383,0.0333,0.0294,0.0263,0.0238$ at
$n=8,12,\dots,40$, so it decays like $1/n$. This is also what rules out a fixed fraction of the
ceiling alone, and that we know only from data. The products $n\cdot\text{ratio}$ are
$0.80,1.08,1.04,0.91,0.92,0.93,0.94,0.95,0.95$. The correct statement compares $\ell'$ with the smaller of
the two bounds.

\begin{conjecture}[Conjecture]\label{conj:resid}
Let $k\ge2$. There is $c=c(k)>0$, depending on $k$ alone, such that for every $n$ with $2k<n$, every
$M$ with $|M|=k$, and every nonempty set $W$ of at most $k-1$ integers each at least $n$,
\[
  \ell_W\ \ge\ c\,\min\Bigl(\ell(n,M),\ \frac{n-2}{n\,\max W}\Bigr),
\]
where $\ell_W$ is the length of the longest component of
$U(n,M)\setminus\bigcup_{w\in W}B_w$.
\end{conjecture}

The quantification over sets $W$, rather than over a single $v_1$, is what the recursion needs: at
step $j$ the region already has $j$ punctures, and a statement about one puncture does not reach it.
The ceiling is governed by $\max W$, since a connected subset of the complement of
$\bigcup_{w\in W}B_w$ lies in the complement of each $B_w$ and Proposition~\ref{prop:ceiling} applies
to the largest.

Granting it, every inserted speed is bounded. Write $\ell_j$ for the longest component after $j$
punctures. Theorem~\ref{thm:cover} applied to that region with the remaining $k-j$ speeds gives
$v_{j+1}\le 2(k-j)/(\ell_j(n-2(k-j)))$, and each branch of the minimum is bounded:
\[
  v_{j+1}\ \le\ \frac{2(k-j)\,n\,\max_{i\le j}v_i}{c\,(n-2)\,(n-2k+2j)}
  \qquad\text{or}\qquad
  v_{j+1}\ \le\ \frac1c\cdot\frac{2(k-j)}{\ell(n,M)\,\bigl(n-2(k-j)\bigr)} .
\]
The second is $1/c$ times the right-hand side of Theorem~\ref{thm:cover} with $k$ replaced by $k-j$,
so Theorem~\ref{thm:uniform}(i) bounds it by $4(k-j)n^{3}/(c(n-2k+2j))$. The first is bounded by
induction, since $v_1\le 4kn^{3}/(n-2k)$ by Corollary~\ref{cor:main}. Iterating bounds every element
of $R$ and makes the classification a finite computation for each $n$, answering for $k$-speed
modifications a version of \cite[Question~9.3]{zhang}. Theorem~\ref{thm:cover} is stated for
$U(n,M)$, but its proof uses only that the region is covered by the remaining $B_v$ and has a
component of length $\ell$, so it applies verbatim.

The cost is not uniform in $k$. Step $j$ multiplies the bound by
$\frac{2(k-j)\,n}{c\,(n-2)(n-2k+2j)}$, which is about $2(k-j)/(cn)$ when $n\gg k$, and about
$(k-j)/(cj)$ when $n-2k=O(1)$. The composed product $\prod_{j=1}^{k-1}$ of these is bounded by a
constant times $c^{-(k-1)}$: numerically, at $n-2k=1$ and $c=0.2$, the product equals
$0.907,0.698,0.516,0.373$ times $c^{-(k-1)}$ at $k=5,10,20,40$. So the bound on $\max R$ is
exponential in $k$ through $c^{-(k-1)}$, with no additional factorial.

Evidence, in both regimes and at both puncture depths. With one puncture, over all $M$ with $|M|=2$
and all $v_1$ from $n$ to $8n$, the ratio $\ell_W\big/\min(\ell(n,M),(n-2)/(nv_1))$ has minimum
$0.207341$ over $21\le n\le44$, attained at $n=32$, $M=\{18,30\}$, $v_1=36$, and minimum $0.21875$
over $8\le n\le20$, attained at $n=8$, $M=\{3,4\}$, $v_1=8$. With two punctures, exhaustively over all $M$ with $|M|=3$, all
$9\le n\le22$ and all pairs $v_1<v_2$ in $[n,3n]$ ($5\,241\,348$ configurations), the minimum is
$3/22=0.136364$, at $n=13$, $M=\{5,6,7\}$, $W=\{13,14\}$. In no instance did the region lose all positive-length components. The constant falls
with puncture depth, consistent with the degradation just described.

Positivity alone is weaker than it looks: for $|M|=k\ge2$ the set $([n-1]\setminus M)\cup W$ has at
most $n-2$ speeds, so the Lonely Runner Conjecture predicts a floor strictly above $1/n$ for it, and
strictness forces $\ell_W>0$. So the positive-length half of Conjecture~\ref{conj:resid} is a
consequence of the conjecture this paper is about; only the quantitative half is extra.

The restriction $|W|\le k-1$ is what does the work, not $k\ge2$ as such. For $k=1$ the residual can
collapse to a finite set of isolated points: for $n\equiv2\pmod 6$, $M=\{n-2\}$ and $v_1=2(n-2)$ it
is nonempty and has no component of positive length, being $\{1/8,3/8,5/8,7/8\}$ at $n=8$. The
corresponding modifications $([n-1]\setminus\{n-2\})\cup\{2(n-2)\}$ are tight, verified for
$n=8,14,20,26,32,38,44,50$, and are members of the Goddyn--Wong families.

In the cases we examined with $k\ge2$, the surviving component sits at a base point $a/q$ with
$q\in M\cup\{n\}$ and $\gcd(a,q)=1$, displaced when $2q\le n-1$ by the presence of the speed $2q$.
The number of such base points that can carry a component of the maximal length is up to $2k$, not
$k+1$: each $b\in M$ coprime to $n$ contributes its own tight point $a/n$ with $a\equiv b^{-1}\pmod n$, as
well as the points $a/b$. What is missing is an arithmetic criterion showing
that $k$ inserted speeds cannot block all of them at once.

\begin{thebibliography}{99}
\bibitem{bhk} T.~Bohman, R.~Holzman and D.~Kleitman, \emph{Six lonely runners}, Electron. J. Combin. \textbf{8} (2001), no.~2, \#R3.
\bibitem{bp} T.~Bohman and F.~Peng, \emph{Coprime mappings and lonely runners}, Mathematika \textbf{68} (2022), 784--804.
\bibitem{bs} J.~Barajas and O.~Serra, \emph{The lonely runner with seven runners}, Electron. J. Combin. \textbf{15} (2008), \#R48.
\bibitem{cusick74} T.~W.~Cusick, \emph{View-obstruction problems in $n$-dimensional geometry}, J. Combin. Theory Ser. A \textbf{16} (1974), 1--11.
\bibitem{cp} T.~W.~Cusick and C.~Pomerance, \emph{View-obstruction problems III}, J. Number Theory \textbf{19} (1984), 131--139.
\bibitem{gw} L.~Goddyn and E.~B.~Wong, \emph{Tight instances of the lonely runner}, Integers \textbf{6} (2006), \#A38.
\bibitem{iwaniec} H.~Iwaniec, \emph{On the problem of Jacobsthal}, Demonstratio Math. \textbf{11} (1978), 225--231.
\bibitem{kravitz} N.~Kravitz, \emph{Barely lonely runners and very lonely runners: a refined approach to the lonely runner problem}, Comb. Theory \textbf{1} (2021), \#P17.
\bibitem{fansun} H.~T.~Fan and A.~Sun, \emph{Amending the lonely runner spectrum conjecture}, arXiv:2306.10417.
\bibitem{ps} G.~Perarnau and O.~Serra, \emph{The Lonely Runner Conjecture turns 60}, Comput. Sci. Rev. \textbf{58} (2025), 100798.
\bibitem{pomerance} C.~Pomerance, \emph{Coprime matchings}, Integers \textbf{22} (2022), \#A29.
\bibitem{rosenfeld} M.~Rosenfeld, \emph{The lonely runner conjecture holds for eight runners}, arXiv:2509.14111.
\bibitem{rs} J.~B.~Rosser and L.~Schoenfeld, \emph{Approximate formulas for some functions of prime numbers}, Illinois J. Math. \textbf{6} (1962), 64--94.
\bibitem{st} T.~Sungkawichai and T.~Trakulthongchai, \emph{Eleven, twelve, and thirteen lonely runners}, arXiv:2604.23906.
\bibitem{tao} T.~Tao, \emph{Some remarks on the lonely runner conjecture}, Contrib. Discrete Math. \textbf{13} (2018), 1--31.
\bibitem{trak} T.~Trakulthongchai, \emph{Nine and ten lonely runners}, Electron. J. Combin. \textbf{33} (2026), \#P2.46.
\bibitem{zhang} Y.~Zhang, \emph{Tight instances of the Lonely Runner Conjecture: complete classification of one-entry modifications, a new infinite family, and the growth bound}, arXiv:2608.13599v2 (31 August 2026).
\end{thebibliography}

\end{document}
